%\documentclass[anon,12pt]{colt2016} % Anonymized submission
\documentclass{colt2016} % Include author names

% The following packages will be automatically loaded:
% amsmath, amssymb, natbib, graphicx, url, algorithm2e

\title[The low-rank approach for certain semidefinite programs]{On the low-rank approach for semidefinite programs arising in synchronization and community detection}

%\title[The low rank approach for certain Semidefinite Programs]%
%{Lift me up but not too high:\\On the low rank approach for semidefinite programs arising in synchronization and community detection}

\usepackage{times}
 % Use \Name{Author Name} to specify the name.
 % If the surname contains spaces, enclose the surname
 % in braces, e.g. \Name{John {Smith Jones}} similarly
 % if the name has a "von" part, e.g \Name{Jane {de Winter}}.
 % If the first letter in the forenames is a diacritic
 % enclose the diacritic in braces, e.g. \Name{{\'E}louise Smith}

 % Two authors with the same address
  % \coltauthor{\Name{Author Name1} \Email{abc@sample.com}\and
  %  \Name{Author Name2} \Email{xyz@sample.com}\\
  %  \addr Address}

 % Three or more authors with the same address:
 % \coltauthor{\Name{Author Name1} \Email{an1@sample.com}\\
 %  \Name{Author Name2} \Email{an2@sample.com}\\
 %  \Name{Author Name3} \Email{an3@sample.com}\\
 %  \addr Address}


 % Authors with different addresses:
 \coltauthor{\Name{Afonso S. Bandeira} \Email{bandeira@mit.edu}\\
 \addr Department of Mathematics, Massachusetts Institute of Technology, Cambridge, Massachusetts, USA.
 \AND
 \Name{Nicolas Boumal} \Email{nboumal@math.princeton.edu}\\
 \addr Mathematics Department, Princeton University, Princeton, New Jersey, USA.
  \AND
 \Name{Vladislav Voroninski} \Email{vvlad@math.mit.edu}\\
 \addr Department of Mathematics, Massachusetts Institute of Technology, Cambridge, Massachusetts, USA.
 }






\newcommand{\ER}{{Erd\H{o}s-R\'enyi }}
\newcommand{\argmin}[1]{\underset{#1}{\operatorname{argmin}}}
\newcommand{\argmax}[1]{\underset{#1}{\operatorname{argmax}}}
\newcommand{\transpose}{^\top\! }
\newcommand{\eig}{{\mathrm{eig}}}
\newcommand{\mle}{{\mathrm{MLE}}}
\newcommand{\mlez}{{\mathrm{MLE0}}}
\newcommand{\inner}[2]{\left\langle{#1},{#2}\right\rangle}
\newcommand{\innersmall}[2]{\langle{#1},{#2}\rangle}
\newcommand{\innerbig}[2]{\big\langle{#1},{#2}\big\rangle}
\newcommand{\MSE}{\mathrm{MSE}}
\newcommand{\Egood}{E_\textrm{good}}
\newcommand{\Ebad}{E_\textrm{bad}}
\newcommand{\trace}{\mathrm{Tr}}
\newcommand{\Trace}{\mathrm{Tr}}
\newcommand{\spann}{\mathrm{span}}
%\newcommand{\skeww}[1]{\left\lceil #1 \right\rfloor}
\newcommand{\skeww}[1]{\operatorname{skew}\!\left( #1 \right)}
\newcommand{\symmop}{\operatorname{sym}}
\newcommand{\symm}[1]{\symmop\!\left( #1 \right)}
\newcommand{\symmsmall}[1]{\symmop( #1 )}
%\newcommand{\vecc}{\mathrm{vec}}
\newcommand{\veccc}[1]{\vecc\left({#1}\right)}
\newcommand{\sbdop}{\operatorname{symblockdiag}}
\newcommand{\sbd}[1]{\sbdop\!\left({#1}\right)}
\newcommand{\sbdsmall}[1]{\sbdop({#1})}
\newcommand{\RMSE}{\mathrm{RMSE}}
\newcommand{\Stdpm}{\St(d, p)^m}
\newcommand{\Proj}{\mathrm{Proj}}
\newcommand{\ProjH}{\mathrm{Proj}^h}
\newcommand{\ProjV}{\mathrm{Proj}^v}
\newcommand{\Exp}{\mathrm{Exp}}
\newcommand{\Log}{\mathrm{Log}}
\newcommand{\expect}{\mathbb{E}}
\newcommand{\expectt}[1]{\mathbb{E}\left\{{#1}\right\}}
%\newcommand{\R}{\mathrm{R}}
\newcommand{\Retr}{\mathrm{Retraction}}
\newcommand{\Trans}{\mathrm{Transport}}
\newcommand{\Transbis}[2]{\mathrm{Transport}_{{#1}}({#2})}
\newcommand{\polar}{\mathrm{polar}}
\newcommand{\qf}{\mathrm{qf}}
\newcommand{\Gr}{\mathrm{Gr}}
\newcommand{\St}{\mathrm{St}}
\newcommand{\GLr}{{\mathbb{R}^{r\times r}_*}}
\newcommand{\Precon}{\mathrm{Precon}}
\newcommand{\T}{\mathrm{T}}
\newcommand{\HH}{\mathrm{H}}
\newcommand{\VV}{\mathrm{V}}
\newcommand{\N}{\mathcal{N}}
\newcommand{\M}{\mathcal{M}}
\newcommand{\barM}{\overline{\mathcal{M}}}
\newcommand{\barnabla}{\overline{\nabla}}
\newcommand{\barf}{{\overline{f}}}
\newcommand{\barx}{{\overline{x}}}
\newcommand{\bary}{{\overline{y}}}
\newcommand{\barX}{{\overline{X}}}
\newcommand{\barY}{{\overline{Y}}}
\newcommand{\barg}{{\overline{g}}}
\newcommand{\barxi}{{\overline{\xi}}}
\newcommand{\bareta}{{\overline{\eta}}}
\newcommand{\p}{\mathcal{P}}
\newcommand{\SOn}{{\mathrm{SO}(n)}}
\newcommand{\On}{{\mathrm{O}(n)}}
\newcommand{\Od}{{\mathrm{O}(d)}}
\newcommand{\Op}{{\mathrm{O}(p)}}
\newcommand{\son}{{\mathfrak{so}(n)}}
\newcommand{\SOt}{{\mathrm{SO}(3)}}
\newcommand{\sot}{{\mathfrak{so}(3)}}
\newcommand{\SOtwo}{{\mathrm{SO}(2)}}
\newcommand{\sotwo}{{\mathfrak{so}(2)}}
\newcommand{\So}{{\mathbb{S}^{1}}}
\newcommand{\Stwo}{{\mathbb{S}^{2}}}
\newcommand{\Sp}{{\mathbb{S}^{2}}}
\newcommand{\Rtt}{{\mathbb{R}^{3\times 3}}}
\newcommand{\Snn}{{\mathbb{S}^{n\times n}}}
\newcommand{\Smm}{{\mathbb{S}^{m\times m}}}
\newcommand{\Smdmd}{{\mathbb{S}^{md\times md}}}
\newcommand{\Spp}{{\mathbb{S}^{p\times p}}}
\newcommand{\Sdd}{{\mathbb{S}^{d\times d}}}
\newcommand{\Srr}{{\mathbb{S}^{r\times r}}}
\newcommand{\Rnn}{{\mathbb{R}^{n\times n}}}
\newcommand{\Rnp}{{\mathbb{R}^{n\times p}}}
\newcommand{\Rdp}{{\mathbb{R}^{d\times p}}}
\newcommand{\Rmn}{{\mathbb{R}^{m\times n}}}
\newcommand{\Rnm}{\mathbb{R}^{n\times m}}
\newcommand{\Rnr}{\mathbb{R}^{n\times r}}
\newcommand{\Rrr}{\mathbb{R}^{r\times r}}
\newcommand{\Rpp}{\mathbb{R}^{p\times p}}
\newcommand{\Rmm}{{\mathbb{R}^{m\times m}}}
\newcommand{\Rmr}{\mathbb{R}^{m\times r}}
\newcommand{\Rrn}{{\mathbb{R}^{r\times n}}}
\newcommand{\Rd}{{\mathbb{R}^{d}}}
\newcommand{\Rp}{{\mathbb{R}^{p}}}
\newcommand{\Rk}{{\mathbb{R}^{k}}}
\newcommand{\Rdd}{{\mathbb{R}^{d\times d}}}
\newcommand{\RNN}{{\mathbb{R}^{N\times N}}}
\newcommand{\reals}{{\mathbb{R}}}
\newcommand{\complex}{{\mathbb{C}}}
\newcommand{\Rn}{{\mathbb{R}^n}}
\newcommand{\Rm}{{\mathbb{R}^m}}
\newcommand{\CN}{{\mathbb{C}^{\,N}}}
\newcommand{\CNN}{{\mathbb{C}^{\ N\times N}}}
\newcommand{\CnotR}{{\mathbb{C}\ \backslash\mathbb{R}}}
\newcommand{\grad}{\mathrm{grad}}
\newcommand{\Hess}{\mathrm{Hess}}
\newcommand{\vecc}{\mathrm{vec}}
\newcommand{\sign}{\mathrm{sign}}
\newcommand{\diag}{\mathrm{diag}}
\newcommand{\D}{\mathrm{D}}
\newcommand{\dt}{\mathrm{d}t}
\newcommand{\ds}{\mathrm{d}s}
\newcommand{\calA}{\mathcal{A}}
\newcommand{\calN}{\mathcal{N}}
\newcommand{\calM}{\mathcal{M}}
\newcommand{\calC}{\mathcal{C}}
\newcommand{\calE}{\mathcal{E}}
\newcommand{\calF}{\mathcal{F}}
\newcommand{\calL}{\mathcal{L}}
\newcommand{\calU}{\mathcal{U}}
\newcommand{\calH}{\mathcal{H}}
\newcommand{\calO}{\mathcal{O}}
\newcommand{\Id}{\operatorname{Id}} % need to be distinguishable from identity matrix
\newcommand{\rank}{\operatorname{rank}}
\newcommand{\nulll}{\operatorname{null}}
\newcommand{\col}{\operatorname{col}}
\newcommand{\Kmax}{K_{\mathrm{max}}}

\newcommand{\1}{\mathbf{1}}

\newcommand{\ZZ}{\mathbb{Z}}
\newcommand{\RR}{\mathbb{R}}
\newcommand{\EE}{\mathbb{E}}
\newcommand{\GGG}{\mathcal{G}}
\newcommand{\YYY}{\mathcal{Y}}
\newcommand{\NNN}{\mathcal{N}}
\newcommand{\eps}{\varepsilon}

\newcommand{\frobnormbig}[1]{\big\|{#1}\big\|_\mathrm{F}}
%\newcommand{\frobnorm}[1]{\left\|{#1}\right\|_\mathrm{F}}
\newcommand{\opnormsmall}[1]{\|{#1}\|_\mathrm{op}}
\newcommand{\frobnormsmall}[1]{\|{#1}\|_\mathrm{F}}
\newcommand{\smallfrobnorm}[1]{\|{#1}\|_\mathrm{F}}
\newcommand{\norm}[1]{\left\|{#1}\right\|}
\newcommand{\sqnorm}[1]{\left\|{#1}\right\|^2}
%\newcommand{\sqfrobnorm}[1]{\frobnorm{#1}^2}
\newcommand{\sqfrobnormsmall}[1]{\frobnormsmall{#1}^2}
\newcommand{\sqfrobnormbig}[1]{\frobnormbig{#1}^2}

\newcommand{\frobnorm}[2][F]{\left\|{#2}\right\|_\mathrm{#1}}
\newcommand{\sqfrobnorm}[2][F]{\frobnorm[#1]{#2}^2}

\newcommand{\dd}[1]{\frac{\mathrm{d}}{\mathrm{d}#1}}
\newcommand{\dmu}{\mathrm{d}\mu}
\newcommand{\pdd}[1]{\frac{\partial}{\partial #1}}
\newcommand{\TODO}[1]{{\color{red}{[#1]}}}
%\newcommand{\TODO}[1]{}
\newcommand{\dblquote}[1]{``#1''}
\newcommand{\floor}[1]{\lfloor #1 \rfloor}
\newcommand{\uniform}{\mathrm{Uni}}
\newcommand{\langevin}{\mathrm{Lang}}
\newcommand{\langout}{\mathrm{LangUni}}
\newcommand{\Zeq}{{Z_\mathrm{eq}}}
\newcommand{\feq}{{f_\mathrm{eq}}}
\newcommand{\XM}{\mathfrak{X}(\M)}
\newcommand{\FM}{\mathfrak{F}(\M)}
\newcommand{\length}{\operatorname{length}}
\newcommand{\bfR}{\mathbf{R}}
\newcommand{\bfxi}{{\bm{\xi}}}
\newcommand{\bfX}{{\mathbf{X}}}
\newcommand{\bfY}{{\mathbf{Y}}}
\newcommand{\bfeta}{{\bm{\eta}}}
\newcommand{\bfF}{\mathbf{F}}
\newcommand{\bfOmega}{\boldsymbol{\Omega}}
\newcommand{\bftheta}{{\bm{\theta}}}
\newcommand{\lambdamax}{\lambda_\mathrm{max}}
\newcommand{\lambdamin}{\lambda_\mathrm{min}}


\newcommand{\ddiag}{\mathrm{ddiag}}
\newcommand{\dist}{\mathrm{dist}}
\newcommand{\embdist}{\mathrm{embdist}}
\newcommand{\var}{\mathrm{var}}
\newcommand{\Cov}{\mathbf{C}}
\newcommand{\Covm}{C}
\newcommand{\FIM}{\mathbf{F}}
\newcommand{\FIMm}{F}
\newcommand{\Rmle}{{\hat\bfR_\mle}}
%\newcommand{\diag}{\mathbf{\mathrm{diag}}}
\newcommand{\Rmoptensor}{\mathbf{{R_\mathrm{m}}}}
\newcommand{\barRmoptensor}{\mathbf{{\bar{R}_\mathrm{m}}}}
\newcommand{\Rmop}{R_\mathrm{m}}
\newcommand{\barRmop}{\bar{R}_\mathrm{m}}
\newcommand{\Rcurv}{\mathcal{R}}
\newcommand{\barRcurv}{\bar{\mathcal{R}}}

\newcommand{\ith}{$i$th }
\newcommand{\jth}{$j$th }
\newcommand{\kth}{$k$th }
\newcommand{\ellth}{$\ell$th }

\newcommand{\expp}[1]{{\exp\!\big({#1}\big)}}


\newtheorem{assumption}{Assumption} %[section]

\newcommand{\defeq}{\mathrel{\mathop:}=}
\newcommand{\eqdef}{\mathrel{\mathop=}:}








\begin{document}

\maketitle

\begin{abstract}
To address difficult optimization problems, convex relaxations based on semidefinite programming are now common place in many fields. Although solvable in polynomial time, large semidefinite programs tend to be computationally challenging. Over a decade ago, exploiting the fact that in many applications of interest the desired solutions are low rank, Burer and Monteiro proposed a heuristic to solve such semidefinite programs by restricting the search space to low-rank matrices. The accompanying theory does not explain the extent of the empirical success. We focus on Synchronization and Community Detection problems and provide theoretical guarantees shedding light on the remarkable efficiency of this heuristic. 
\end{abstract}

\begin{keywords}
Semidefinite Programming, Burer--Monteiro heuristic, SDPLR, Synchronization, Community Detection
\end{keywords}


\section{Introduction}

Estimation problems in many fields, including signal processing, statistics and machine learning, are formulated as intractable optimization problems. A now popular technique to attempt solving some of these problems is to replace the difficult optimization problem by a surrogate tractable convex problem and take advantage of existing machinery for convex optimization. In many instances, the surrogate problem is obtained by considering a larger, convex feasible space, and thus it is often called a \textit{convex relaxation}. Relaxations based on semidefinite programming are among the most popular. %In this paper, we study a class of semidefinite relaxations which can be solved through a low-dimensional (but nonconvex) optimization problem instead.


We will focus mostly on a certain class of problems, namely, \emph{synchronization problems} on graphs (\cite{singer2010angular,bandeira2015nonuniquegames}). Synchronization problems consist in estimating $n$ group elements $g_1,\dots,g_n$ from information about their offsets $g_i^{}g_j^{-1}$. A simple instance of this class is $\ZZ_2$-Synchronization (\cite{abbe2014decoding}), corresponding to the group on two elements. In that case, the task is to estimate $n$ bits $x_1,\dots,x_n\in\{\pm1\}$ from noisy measurements of $x_ix_j$.


The problem of Community Detection in the Binary Stochastic Block Model (SBM) can also be regarded as an instance of Synchronization over $\ZZ_2$. SBM on two communities (also known as \emph{planted partition}) is a model of a random graph $\GGG(n,p,q)$ on $n=2m$ nodes split evenly in two communities, identified by $\left\{\pm1\right\}$ node labels. Edges for this graph are drawn randomly and independently, where each pair of nodes in the same community gets connected with probability $p$ and in different communities with probability $q$. The task is to estimate the partition, given an observation of the graph. This fascinating problem was the target of much study in the last few years, including identification of remarkable phase transitions on the values of $p$ and $q$ and possibility of recovering the unknown labels (\cite{Decelle_SBM,Mossel_SBM1,Mossel_SBM2,Massoulie_SBM,abbe2014exact,Mossel_SBM3_exact}).


While computing the maximum likelihood estimator (MLE) for either of the problems above is known to be computationally intractable, several heuristics have been proposed and studied. A particularly successful approach is based on Semidefinite Programming. Following the seminal work of \cite{goemans1995maxcut} in the context of the Max-Cut problem, the maximum likelihood estimation problem (originally an intractable optimization problem in $n$ node variables) is relaxed to a semidefinite program (SDP) (a tractable problem on $n^2$ variables). Importantly, the solution to the SDP is only guaranteed to correspond to the solution of the original problem if it has rank~1. Enforcing this constraint explicitly would render the problem intractable. Remarkably, in some regimes, it is known that the solution of this semidefinite program is naturally of rank 1, and that furthermore this solution allows to identify the true labels (\cite{abbe2014decoding,bandeira2014tightness,abbe2014exact,bandeira2015laplacian,Hajek_et_al_SBM_SDP}).


While SDP's are known to be solvable up to arbitrary precision in polynomial time (\cite{nesterov2004introductory}), the increase in dimension due to the lifting technique (the relaxation) and the semidefiniteness constraint render solving them rather slow in practice. This is mainly because intermediate iterates involve matrix decompositions of dense, full-rank matrices of size $n$. On the other hand, in certain regimes, the optimal solution is expected to have low rank. Indeed, in the problems studied here, the solution being rank 1 coincides precisely with it corresponding to the desirable MLE. It is then natural to attempt to reach the solution via a sequence of similarly simple objects instead.

The most popular and successful such low-rank approach is SDPLR, as proposed in~\citep{sdplr,burer2005local} (also in a particular form in \citep{burer2002rank}). In a nutshell, the idea is to restrict the search space to matrices of bounded rank. While, after adding this rank constraint, the optimization problem is no longer convex (and it is hence unclear whether it is tractable or not),
%or the method guaranteed to produce meaningful solutions,
this approach is empirically successful for a variety of instances~\citep{sdplr,journee2010low,bandeira2014tightness,boumal2015staircase}.

The original SDPLR papers come with general theory supporting the fact that relaxing the rank up to about $\sqrt{2n}$ might work well. Yet, in practice, for well-behaved SDP's which admit a solution of rank 1, it is often seen that relaxing the rank merely to 2 works fantastically. To date, there was no satisfactory explanation for this nonconvex success. This paper provides the first such guarantee, in the context of synchronization over $\ZZ_2$ and community detection. More precisely, we show that, in certain noise regimes, there are no spurious second-order critical points for the
rank-2 constrained problem, while in other (more inclusive) regimes we show that all such points need to correlate non-trivially with the ground truth. Second-order critical points (that is, points which satisfy second-order necessary optimality conditions) can be computed in our context~\citep{boumal2016complexity}.

In this paper, we focus on a selection of well-studied problems, keeping in mind that the proposed analysis has the potential to generalize to many problems where semidefinite relaxations are successful yet demanding to solve. A more general setting will be the focus of a future publication.

It is also worth noting that semidefinite relaxations have been observed to work well on problems for which other standard heuristics seem to not perform well. One relevant example is the Multireference Alignment problem~(\cite{Bandeira_Charikar_Singer_Zhu_Alignment}). Furthermore, low-rank semidefinite programming has the potential to yield a posteriori certifiably correct algorithms~(\cite{Bandeira_PCC}) and is known to be robust to certain monotone adversary models~(\cite{Feige_Kilian_bisection_01,Moitra_SBMadversary}).



\subsection*{Notation}

Given a matrix $M$ and its vector of singular values $\nu$, we write $\|M\| = \|\nu\|_\infty$ for its operator norm (largest singular value), $\|M\|_F = \|\nu\|_2$ for its Frobenius norm and $\|M\|_* = \|\nu\|_1$ for its nuclear norm (sum of singular values). The operator $\ddiag \colon \Rnn \to \Rnn$ sets all off-diagonal entries of a matrix to zero. The Hadamard or entry-wise product is written $\circ$.


\section{The $\ZZ_2$ Synchronization problem}

The goal of $\ZZ_2$ Synchronization is to estimate labels $z\in\{\pm 1\}^n$ from noisy pairwise measurements $Y_{ij} = z_iz_j + \sigma W_{ij}$, where $W_{i>j} \stackrel{i.i.d.}{\sim} \NNN(0,1)$ and $W_{ij}=W_{ji}$, $W_{ii} = 0$. In matrix notation,
\begin{equation}\label{eq:noisemodel:sigma}
 Y = zz^T + \sigma W.
\end{equation}
Because only $zz^T$ is measured, the labels can only be estimated up to a global sign flip. One can also consider more realistic noise models on which $Y_{ij}$ also takes binary values (\cite{abbe2014decoding}) but, for the sake of simplicity, we will consider Gaussian noise. Since $\|zz^T\| = n$ and $\|\sigma W\|$ concentrates around a constant multiple of $\sigma \sqrt{n}$ (in operator norm), a natural way to parametrize the signal-to-noise ratio is by taking $\lambda = \frac{\sqrt{n}}{\sigma}$. For this reason, some authors consider the scaled model
\begin{equation}\label{eq:noisemodel:lambda}
 \YYY = \frac{\lambda}{n}Y =  \frac{\lambda}{n}zz^T + \frac1{\sqrt{n}}W.
\end{equation}
The problems of recovering $z$ from $Y$ or $\YYY$ are clearly equivalent. The former choice of notation has been used, for example, in~\citep{bandeira2014tightness,bandeira2015laplacian} and has the advantage of highlighting the fact that the observation is an entry-wise noisy version of the ground truth, while the latter has been used in~\citep{javanmard2015phase} and has the advantage of highlighting the spike model interpretation of the problem and making a more transparent connection with the well-studied spike model in random matrix theory and the BBP transition phenomenon~\citep{BBP_MC_BBP_2005,Feral_Peche_BBPWigner} (this connection had already been made in~\citep{singer2010angular}). For the sake of completeness, we will state our results in both notation choices.

The most natural approach to recovering $z$ is to consider the MLE, which solves
\begin{equation}
\min_{x\in\{\pm 1\}^n} \sum_{i,j=1}^n \left( Y_{ij} - x_ix_j  \right)^2.
\end{equation}
It is readily seen that solutions of this problem are the same as those of
\begin{equation}\label{eq:littleGrothendieck}
\max_{x\in\{\pm 1\}^n} x^TYx,
\end{equation}
which is known to the \emph{NP-hard} in general. In fact, when $Y\succeq 0$ this is known as the little Grothendieck problem~\citep{briet2014grothendieck,bandeira2013approximating}, and, when $Y$ corresponds to the Laplacian of a graph, it corresponds to the Max-Cut problem~\citep{goemans1995maxcut}.

Following the now standard lifting technique (dating back at least to~\cite{goemans1995maxcut} ), we set a new variable $X = xx^T$ and write the equivalent formulation
\begin{equation}\label{SDP:rank1}
\begin{array}{ll}
\max_X & \trace\left( YX \right) \\
\text{s.t.} & X_{ii} = 1, \text{ for } 1\leq i \leq n\\ 
 & X \succeq 0 \\
 & \rank(X) = 1.
\end{array}
\end{equation}
Problems~\eqref{eq:littleGrothendieck} and~\eqref{SDP:rank1} are equivalent.
One arrives at a tractable SDP formulation by dropping the problematic rank constraint:
\begin{equation}\label{SDP:fullrank}
\begin{array}{ll}
\max & \trace\left( YX \right) \\
\text{s.t.} & X_{ii} = 1, \text{ for } 1\leq i \leq n\\ 
 & X \succeq 0.
\end{array}
\end{equation}
By construction, problem \eqref{SDP:fullrank} has the same solutions if $Y$ is replaced with $\YYY$ (but it will have a different optimal value, since $Y$ and $\YYY$ are scaled versions of each other.)

Recall the model~\eqref{eq:noisemodel:sigma} (or equivalently~\eqref{eq:noisemodel:lambda}). It is known that, as long as $\sigma < \sqrt{\frac{n}{2\log n}}$, or equivalently $\lambda > \sqrt{2\log n}$, with high probability, the solution $X$ of~\eqref{SDP:fullrank} is unique and corresponds exactly to the ground truth $X=zz^T$. We refer to this phenomenon as \emph{exact recovery}. It is also known that on the other side of this threshold exact recovery is information-theoretically impossible, showcasing the efficiency of semidefinite relaxations for this type of task~\citep{bandeira2014tightness,bandeira2015laplacian}.

While exact recovery is impossible for constant $\lambda$ (or, equivalently, $\sigma \sim \sqrt{n}$), simply taking the top eigenvector of $\YYY$ is known to produce an estimator that correlates non-trivially with the ground truth, precisely when $\lambda >1$ (\cite{Feral_Peche_BBPWigner,singer2010angular}). This phenomenon is often referred to as the BBP transition (\cite{BBP_MC_BBP_2005}). This motivates the question of whether~\eqref{SDP:fullrank} gives meaningful results for the constant $\lambda$ regime~(\cite{Montanari_SDPdetectionSBM,javanmard2015phase,Guedon_SBMgrothendieck}). In the context of community detection, this question was first addressed by~\cite{Guedon_SBMgrothendieck}. In~\citep{Montanari_SDPdetectionSBM}, a phase transition is shown to exist at $\lambda=1$: similarly to what happens with the top eigenvalue of $\YYY$ (\cite{Feral_Peche_BBPWigner}). For $\lambda<1$, the value of~\eqref{SDP:fullrank} with cost $\YYY$ converges (in probability) to $2$, and for $\lambda>1$ its limit is strictly larger than $2$. \cite{javanmard2015phase} give fascinating predictions, based on non-rigorous statistical mechanics tools, for the behavior of both~\eqref{SDP:fullrank} and~\eqref{eq:littleGrothendieck} as a function of $\lambda$.

Another, related, type of synchronization problem is \emph{phase synchronization}, where one estimates $z\in\complex^n$ with $|z_i|=1$ for all $i$. This is a direct complex analogue of $\ZZ_2$ synchronization. The SDP relaxation works similarly in the complex field. See~\citep{singer2010angular} for a first analysis of the SDP relaxation, see~\citep{bandeira2014tightness} for a proof of tightness of the SDP relaxation in a similar regime as here and, still in the same regime, see~\citep{boumal2016nonconvexphase} for a proof that the nonconvex problem has no spurious local optima either. A key difference with the present paper is that phase synchronization is a continuous problem, whereas $\ZZ_2$ synchronization is discrete. This explains why, in the present setting, the constraint $\rank(X)=1$ has to be relaxed at least to $\rank(X)=2$ (to connect the search space), whereas in the former paper, the rank is not relaxed at all.

\subsection{The Burer--Monteiro Approach}

In transiting from~\eqref{SDP:rank1} to~\eqref{SDP:fullrank}, one effectively replaces the constraint $\rank(X) \leq 1$ by the vacuous constraint $\rank(X) \leq n$, thus going from a combinatorial problem to a tractable but high-dimensional SDP. One of the insights of~\cite{sdplr,burer2005local} is that relaxing the rank only partially, as $\rank(X)\leq p$ for variable $p$, gives access to a family of low-dimensional (but nonconvex) nonlinear optimization problems, which can be put to good use to understand both~\eqref{SDP:rank1} and~\eqref{SDP:fullrank}.

In general, given an SDP of the form
\begin{align}
	\max_{X} \Trace(YX) \quad \textrm{ s.t. } \quad \calA(X) = b, X \succeq 0,
\end{align}
where $\calA$ is a linear operator from the symmetric matrices of size $n$ to $\Rm$,
the SDPLR algorithm (sometimes referred to as the Burer--Monteiro approach) consists in parameterizing $X$ as $X = QQ^T$, where $Q$ lives in $\Rnp$. In so doing, $X \succeq 0$ is naturally enforced. This yields the following nonlinear optimization problem:
\begin{align}
\max_{Q\in\Rnp} \Trace(Q^TYQ) \quad \textrm{ s.t. } \quad \calA(QQ^T) = b,
\end{align}
where both the cost and the constraints are quadratic in $Q$. This is typically nonconvex. Both problems are equivalent, up to the additional constraint $\rank(X) \leq p$ forced in the nonlinear program. Of course, if the SDP admits a solution of rank at most $p$, then the two problems attain the same optimal value. When the search space is compact, this is known to happen as soon as $p(p+1)/2 \geq m$ ($m$ is the number of constraints), as per general results put forth in~\citep{shapiro1982rank,barvinok1995problems,pataki1998rank}. \cite{burer2005local} show that rank deficient local optima of the nonlinear program are globally optimal.\footnote{This renders these methods potentially a posteriori certifiable (\cite{Bandeira_PCC}).} See also Lemma~\ref{lem:rankdeficientglobalopt}.

The latter observations motivate the algorithm SDPLR~\citep{sdplr}, where the nonlinear program is tackled in lieu of the SDP, using classical nonlinear optimization algorithms such as the augmented Lagrangian method. The above discussion suggests setting $p \approx \sqrt{2m}$. If the local method converges to a rank-deficient local optimum (which is often the case in practice but is not satisfactorily explained in theory), then we have found a global optimum. The advantage (compared to the SDP) is that the search space has lower dimension.

%The theory in~\citep{burer2005local} states that, for such $p$, suboptimal local optima (necessarily full rank) must have certain strong geometric properties (see also~\citep{boumal2015staircase} for a discussion of such points in the deterministic setting). It has been said in the literature that points satisfying these conditions (thus, potential culprits for SDPLR) are rare, but to the best of our knowledge, there is no supporting theory for that statement. \TODO{UPDATE WITH RESPECT TO NEWEST PAPER}

In practice, if we are optimistic about our chances, we can set $p$ as low as $p=2$ for example~(\cite{burer2002rank}). If the SDP has a solution of rank 1, then the corresponding nonlinear program has local optima (actually, global optima) of rank 1 as well, thus being rank deficient, which allows certifying their optimality. In practice, this is seen to work remarkably well for the problems considered in this paper (in certain noise regimes), as had already been observed in other papers as well. It is not obvious why this is so, as the rank restriction could (in general) spawn a large number of spurious local optima.

Here, for the first time, we provide a proof that, at $p=2$ and in favorable regimes, all local optima are global optima, which explains why local methods behave appropriately. In fact, we even show that saddle points can always be escaped using solely second-order derivatives.

Local algorithms such as the augmented Lagrangian method are useful to tackle generic constrained nonlinear programs. The SDP's under scrutiny, though, have rather special structure. In particular, they are such that the search space in $Q$ is a smooth manifold already for $p\geq 2$. As originally advocated by~\cite{journee2010low}, this suggests solving the nonlinear program using techniques from optimization on manifolds~\citep{AMS08}. The latter are indeed better suited to fully leverage the special geometry of these problems. See also~\citep{wen2013orthogonality,boumal2015staircase} for further numerical and theoretical investigation.
%An immediate theoretical advantage of recognizing this special structure is simplified expressions for the necessary optimality conditions of the nonlinear program, as readily noted in~\citep{burer2002rank}. See Lemma
See~\citep{boumal2016complexity} for a Riemannian version of the trust-region method which globally converges to global optima owing to the special properties of our problem, with global rates of convergence.

As described above, here we consider the rank constrained problem as follows (\cite{burer2002rank}):
\begin{equation}\label{rank2}
\begin{array}{ll}
\max & \trace\left( YX \right) \\
\text{s.t.} & X_{ii} = 1, \text{ for } 1\leq i \leq n\\ 
 & X \succeq 0 \\
 & \rank(X) \leq 2. 
\end{array}
\end{equation}
By taking $X = QQ^T$ with $Q\in\RR^{n\times 2}$, one can reformulate~\eqref{rank2} as:
\begin{equation}\label{rank2BM}
\begin{array}{ll}
\max & \trace\left( Q^TYQ \right) \\
 & \text{s.t.} \left\|Q_{i:}\right\|^2 = 1, \text{ for } 1\leq i \leq n \\
 & Q \in \RR^{n \times 2},  
\end{array}
\end{equation}
where $Q_{i:}$ denotes the $i$th row of $Q$. Note the geometry here: the search space is a product of circles.

% \TODO{
% Address why \cite[Thm.\,3]{wen2013orthogonality} is not satisfactory: as presented, limited to diagonal constraints ; no characterization of when the inf rank takes on a useful value ; even in the Max Cut case, if cost is rank 1 with 1's on diagonal, the theorem's conclusion is vacuous. Do check the reference they give: ``On the rank of a cograph'', but seems to yield nothing useful.
% }

We now present our main results in the realm of $\ZZ_2$ Synchronization.

\begin{theorem}\label{thm:Z2Synch:partialrecovery}
Consider model~\eqref{eq:noisemodel:lambda} (or equivalently~\eqref{eq:noisemodel:sigma}). If $\lambda > 8$ (or equivalently $\sigma < \frac18\sqrt{n}$), then, with high probability, all second-order critical points $Q$ of~\eqref{rank2BM} correlate non-trivially with the ground truth $z$ in the sense that, for every such $\lambda$, there exists $\eps$ such that
\begin{equation}\label{eq:nontrivialcorrelation_0}
\frac1{n} \left\| Q^Tz\right\|_2 \geq \eps.
\end{equation}
\end{theorem}
\begin{proof}
All the interesting ingredients of the proof are in Section~\ref{section:mainresults}. The claim will follow from Lemma~\ref{lemma:Xcorrelates11T}, noting that $\left\| Q^Tz\right\|_2^2 = \left\langle QQ^T,zz^T \right\rangle$, the behavior of~\eqref{rank2BM} is unchanged by changing the diagonal values of $\YYY$, and the fact that for any $\delta>0$, with high probabily, $\frac1n\mathrm{SDP}(W-\ddiag(W))\leq \|W-\ddiag(W)\| \leq (2+\delta)\sqrt{n}$ (see Lemma~\ref{lemma:Wignermatrices} and eq.~\eqref{SDP:eq:inequality}).
\end{proof}

\begin{remark}[Rounding]
In Theorem~\eqref{thm:Z2Synch:partialrecovery} we ask for non-trivial correlation in the form of~\eqref{eq:nontrivialcorrelation_0}, but it would also be natural to ask for an estimator $\hat{z}\in\RR^n$ (or even $\hat{z}\in\{\pm1\}^n$) that exhibits non trivial correlation with $z$. We note that if we take $x$ and $y$ to be the columns of $Q$ satisfying~\eqref{eq:nontrivialcorrelation_0}, then a random linear combination $g_1x+g_2y$, where $g_1$ and $g_2$ are independent standard Gaussian variables, can be shown to be likely to produce meaningful correlation estimators. Indeed, $\EE\{ \|Qg\|^2 \} = \|Q\|_F^2 = n$ and $\EE\{  \inner{z}{Qg}^2 \} = \|Q^Tz\|_2^2$. We also direct the reader to Section~4 of~\citep{Montanari_SDPdetectionSBM} for techniques to construct meaningful $\{\pm1,0\}^n$ estimators from solutions of the SDP. We note furthermore that if $\eps$ is large in~\eqref{eq:nontrivialcorrelation_0} (as will be the case in Theorem~\ref{lem:Z2Synch:almostexactrecovery}, for example), then constructing such estimators becomes considerably easier.
\end{remark}

\begin{theorem}\label{lem:Z2Synch:almostexactrecovery}
Consider model~\eqref{eq:noisemodel:lambda} (or equivalently~\eqref{eq:noisemodel:sigma}). If $\lambda > 16$ (or equivalently $\sigma < \frac1{16}\sqrt{n}$), then for any $\eps>0$, with high probability,
%for large $n$ (potentially depending on $\eps$),
all second-order critical points $Q$ of~\eqref{rank2BM} satisfy
\begin{equation}\label{eq:nontrivialcorrelation_0}
\frac1{n} \left\| Q^Tz\right\|_2 \geq 1 - (1+\eps)\frac{16}{\lambda} =  1 - (1+\eps)\frac{16\sigma}{\sqrt{n}}.
\end{equation}
\end{theorem}
\begin{proof}
Again, all the interesting ingredients of the proof are in Section~\ref{section:mainresults}, as this will follow directly from Lemma~\ref{lemma:Qcorrelates1} and the considerations in the proof of Theorem~\ref{thm:Z2Synch:partialrecovery}.
\end{proof}

\begin{theorem}\label{lem:Z2Synch:exactrecovery}
Consider model~\eqref{eq:noisemodel:lambda} (or equivalently~\eqref{eq:noisemodel:sigma}). There is a universal constant $C$ such that: If $\lambda \geq Cn^{\frac13}$ (or equivalently $\sigma \leq \frac1{C}n^{\frac16}$) then, with high probability, all second-order critical points $Q$ of~\eqref{rank2BM} are optimal and correspond to the ground truth, meaning $QQ^T = zz^T$.
\end{theorem}
\begin{proof}
	This follows from Theorem~\ref{thm:maxcutrktwo} and the high probability bounds in Lemma~\ref{lemma:Wignermatrices}.
\end{proof}





%\TODO{
%Explain that this does not need the typical ``initialization to get close enough''
%}

%\TODO{
%Refer Nicolas' paper appropriately (phases), and argue that what is studied here has the potential to be super general.
%}

%\TODO{Nicolas: comment about local methods, complexity etc: we explain the behavior ; we don't give an algorithm}

Note that there are no guarantees in general that nonlinear optimization algorithms (such as the augmented Lagrangian method) converge to local optima. Yet, we know that only local optima are stable fixed points for such methods; and that this is true regardless of initialization. Hence, one perspective is that the contribution of this paper is to show that all stable fixed points of local methods---without the need for a good initial guess---are global optima (in the given regime). As we mentioned earlier though, because in our case second-order necessary optimality conditions are sufficient for global optimality, results in~\citep{boumal2016complexity} show that the Riemannian trust-region method converges to global optimizers, regardless of initialization, with global rates of convergence.


\subsection{Community Detection in the Binary Stochastic Block Model}

The Stochastic Block Model (SBM) on two communities is a random graph model
$\GGG(n,p,q)$ on $n=2m$ nodes. The nodes are divided evenly in two
communities, identified by a vector of labels $g\in\{\pm1\}^n$ satisfying $g^T\1=0$.
Edges on this graph are drawn independently at
random. A pair of nodes in the same community is connected by an edge
with probability $p$; nodes in different communities with probability
$q$. We focus on the case $p>q$. This means that the adjacency
matrix $A$ of this graph has the following distribution:
\[
 A_{ij} = \left\{ \begin{array}{ll} 1 & \text{with probability } p \\
0 & \text{with probability } 1-p, \end{array} \right.
\]
if $g_ig_j = 1$ (thus including the diagonal), and
\[
 A_{ij} = \left\{ \begin{array}{ll} 1 & \text{with probability } q \\
0 & \text{with probability } 1-q, \end{array} \right.
\]
if $g_ig_j = -1$. %\TODO{Afonso: specify $A_{ii}$? Is it the first case then?}

The goal of community detection is to recover the labels $g$ from a
realization $G\sim\GGG(n,p,q)$. Roughly half a decade
ago,~\cite{Decelle_SBM} conjectured a remarkable phase transition in
the constant average--degree regime: if $p = \frac{a}n$ and $q =
\frac{b}n$ with $a>b$ constants,~\cite{Decelle_SBM} conjectured that
as long as
\begin{equation}\label{eq:condition:partialrecovery:SBM}
\lambda(a,b) \defeq \frac{a-b}{\sqrt{2(a+b)}} > 1,
\end{equation}
it is possible to construct an estimator that, with high probability,
correlates with the true partition, whereas below this threshold that would be impossible. This conjecture was proven in a series of
papers~\citep{Mossel_SBM1,Mossel_SBM2,Massoulie_SBM}.

Another natural question is to understand when it is possible to
exactly recover $g$ (or $-g$). A similar phase transition was
established in~\citep{abbe2014exact} and~\citep{Mossel_SBM3_exact}: if
$p = \alpha\frac{\log n}n$ and $q = \beta\frac{\log n}n$ with
$\alpha>\beta$ constants, then as long as
\begin{equation}\label{eq:conditionsalphabetaSBMcor}
 \sqrt{\alpha} - \sqrt{\beta} > \sqrt{2},
\end{equation}
the MLE for the partition, solution of
\begin{equation}\label{eq:MLE_SBM}
\begin{array}{ll}
 \max & x^TAx \\
 \text{ s.t. } & x \in \{ \pm 1 \}^n \\
 & x^T \1 = 0,
\end{array}
\end{equation}
coincides exactly with the partition, with high probability; and,
moreover, below this threshold it is information-theoretically
impossible to recover the partition.

\cite{abbe2014exact} studied a semidefinite relaxation
of~\eqref{eq:MLE_SBM} analogous to~\eqref{SDP:fullrank}, and
conjectured that its solution coincides with the ground truth (that is,
the SDP achieves exact recovery) at the information-theoretical
limit~\eqref{eq:conditionsalphabetaSBMcor}. This was confirmed
independently in~\citep{Hajek_et_al_SBM_SDP}
and~\citep{bandeira2015laplacian}. The performance of the semidefinite
relaxation on the constant average--degree regime has also been target
of recent work (\cite{Guedon_SBMgrothendieck,Montanari_SDPdetectionSBM,javanmard2015phase}).
\cite{Guedon_SBMgrothendieck} showed that when the left hand side
of~\eqref{eq:condition:partialrecovery:SBM} is a sufficiently large
constant, the solution of the SDP correlates nontrivially with the
partition. This was improved by~\cite{Montanari_SDPdetectionSBM} who
showed that, for every $\epsilon>0$, and large enough average degree,
$\lambda(a,b) > 1+\epsilon$ suffices. 
%Interestingly, the robusteness of the SDP against monotone
%adversaries guarantee that, for finite average degree,
%$\lambda(a,b) > 1$ is not is not
%sufficient~(\cite{Moitra_SBMadversary}).
\cite{javanmard2015phase} give precise predictions for the
behavior of the SDP in this regime, albeit using non-rigorous
techniques. Interestingly, the results in~\citep{Moitra_SBMadversary}
suggest that there may be values of $a$ and $b$
satisfying~\eqref{eq:condition:partialrecovery:SBM} for which the SDP
does not, with high probability, provide meaningful solutions.


% \begin{equation}\label{SDP:fullrank_SBM}
% \begin{array}{ll}
% \max & \trace\left( AX \right) \\
% \text{s. t.} & X_{ii} = 1, \text{ for } 1\leq i \leq n\\
%  & X \succeq 0.
% \end{array}
% \end{equation}

%The performance of the SDP for partial recovery was looked into
%in~\cite{javanmard2015phase} and~\cite{Montanari_SDPdetectionSBM}
%and~\cite{Guedon_SBMgrothendieck}.

We will now write $A$ in a form close to~\eqref{eq:noisemodel:lambda} to help illustrate how community detection in the SBM is closely related to $\ZZ_2$ Synchronization. We start by noting that 
\[
\EE A = \frac{p+q}2\1\1^T + \frac{p-q}2gg^T.
\]
Since $\1$ is a fixed vector, we can consider $A^\natural$ defined as
\[
A^\natural = A - \frac{p+q}2\1\1^T.
\]
It is readily seen that, because of the balanced partition constraint ($x^T\1=0$), replacing $A$ by $A^\natural$ does not affect the solution of the MLE~\eqref{eq:MLE_SBM}. The benefit is that now, since $A^\natural$ no longer has a bias in the direction of $\1\1^T$, we will remove the extra constraint and focus on trying to understand the efficiency of the simpler program:
\begin{equation}\label{eq:MLE_SBM_nobalancing}
\begin{array}{ll}
\max & x^T A^\natural x \\
\text{s.t.} & x \in \{\pm1\}^n,
\end{array}
\end{equation}
and its natural SDP relaxation
\begin{equation}\label{SDP:fullrank_SBM}
\begin{array}{ll}
\max & \trace\left( A^\natural X \right) \\
\text{s.t.} & X_{ii} = 1, \text{ for } 1\leq i \leq n\\ 
 & X \succeq 0.
\end{array}
\end{equation}
Following the Burer--Monteiro approach, we consider the natural analogue to~\eqref{rank2BM}:
\begin{equation}\label{rank2BM_SBM}
\begin{array}{ll}
\max & \trace\left( Q^TA^\natural Q \right) \\
 & \text{s.t.} \left\|Q_{i:}\right\|^2 = 1, \text{ for } 1\leq i \leq n \\
 & Q \in \RR^{n \times 2}. 
\end{array}
\end{equation}
We further write
\begin{equation}\label{eq:def:E}
A^\natural = \frac{p-q}{2}gg^T + \frac{(p-q)n}{2}E + \frac{(p-q)n}{2}D,
\end{equation}
where $D$ is a diaognal matrix, and $E$ is a zero-diagonal centered random matrix (we choose to normalize by the average degree $\frac{(p-q)n}{2}$ to be compatible with the notation in~\citep{Montanari_SDPdetectionSBM}, whose statements about $E$ we will use). We note also that the diagonal matrix $D$ does not affect the solutions of~\eqref{eq:MLE_SBM},~\eqref{SDP:fullrank_SBM}, or~\eqref{rank2BM_SBM}. This means in particular that, in the constant average--degree regime ($p=\frac{a}n$ and $q=\frac{b}{n}$),
\[
\sqrt{\frac2{a+b}}A^\natural - D = \frac{\lambda(a,b)}n gg^T + E,
\]
for $\lambda(a,b)$ defined in~\eqref{eq:condition:partialrecovery:SBM}. This illustrates the parallelism with~\eqref{eq:noisemodel:lambda}. The main difficulty with this setting is that the spectral norm of $E$ is known not to be bounded (as $n\to\infty$), with high probability (similarly to how adjacency matrices of constant average--degree \ER graphs typically have a fluctuation around their mean whose spectral norm is not bounded, see for example~\citep{Krivelevich_Sudakov_ER,Montanari_SDPdetectionSBM}). For this reason, we will focus on another quantity. 

Following the notation in~\citep{Montanari_SDPdetectionSBM}, given a matrix $M$ we define $\mathrm{SDP}(M)$ as
\begin{equation}\label{def:SDP:eq}
\begin{array}{lll}
\mathrm{SDP}(M)
= &\max & \trace\left( M X \right) \\
& \text{s.t.} & X_{ii} = 1, \text{ for } 1\leq i \leq n\\ 
 & & X \succeq 0,
\end{array}
\end{equation}
and note that, since in the search space $\|X\|_\ast = n$, we have
\begin{equation}\label{SDP:eq:inequality}
\mathrm{SDP}(M) \leq n\|M\|.
\end{equation}
(Use H\"older's inequality: $\inner{M}{X} \leq \|M\|\|X\|_*$ and $\|X\|_* = \trace(X) = n$ since $X\succeq 0$.)

Fortunately, Lemma~\ref{lem:optimality}, which is crucial to establish non-trivial correlation guarantees, does not depend on $\|E\|$, which is not bounded, but rather depends only on $\frac1n\mathrm{SDP}(E)$, which is known to be bounded~\citep{Montanari_SDPdetectionSBM} (see Lemma~\ref{lemma:sbm:SDPE} for details). Indeed, using Lemmas~\ref{lem:optimality} and~\ref{lemma:sbm:SDPE} we immediately get the following guarantee.

\begin{theorem}
 Consider the Stochastic Block Model on two communities described above in the constant average--degree regime $p=\frac{a}n$ and $q=\frac{b}n$. Let $d$ denote the average degree parameter
 \[
  d = \frac{a+b}2.
 \]
Then, for any $\delta>0$ there exists $d_0$ such that if $d>d_0$ and $\lambda(a,b) > 8+\delta$ (see~\eqref{eq:condition:partialrecovery:SBM}) then there exists $\eps>0$ such that: with high probability, all second-order critical points $Q$ of~\eqref{rank2BM_SBM} correlate non-trivially with the ground truth partition $g$ in the following sense.
\begin{equation}\label{eq:nontrivialcorrelation_0}
\frac1{n} \left\| Q^Tg\right\|_2 \geq \eps.
\end{equation}
\end{theorem}
\begin{proof}
The proof is analogous to Theorem~\ref{thm:Z2Synch:partialrecovery}, but based on the bound on $\mathrm{SDP(E)}$ in Lemma~\ref{lemma:sbm:SDPE}.
\end{proof}

Further controlling $\|E\|$ and $\|Eg\|_\infty$ (Lemmas~\ref{lemma:sbm:spectralnormE} and~\ref{lemma:sbm:infnormE}) and using Lemma~\ref{thm:maxcutrktwo} we can also obtain a (suboptimal) exact recovery guarantee. It is useful to define the parameter (analogous to~\eqref{eq:condition:partialrecovery:SBM}):
\[
 \tilde{\lambda}(p,q) \defeq \frac{p-q}{\sqrt{2(p+q)}}\sqrt{n}.
\]
Note that if $p=\frac{a}n$ and $q=\frac{b}n$ then $\tilde{\lambda}(p,q) =  \lambda(a,b)$.
\begin{theorem}
  Consider the Stochastic Block Model on two communities described above. There exists a universal constant $c$ such that, as long as
  \[
  \tilde{\lambda}(p,q) \geq cn^{1/3},
  \]
then, with high probability, all second-order critical points $Q$ of~\eqref{rank2BM_SBM} correspond to the ground truth partition, meaning $QQ^T = gg^T$.
\end{theorem}

\begin{proof}
If $\frac{p-q}{\sqrt{2(p+q)}}\sqrt{n} \geq cn^{1/3}$, then $\sqrt{\frac{n}{2}(p+q)}\geq cn^{1/3}$ (use $p+q \geq p-q$). A fortiori, this implies that $\sqrt{\frac{n}2(p+q)} \gg \log n$.
Together with Lemmas~\ref{lemma:sbm:spectralnormE} and~\ref{lemma:sbm:infnormE} this shows that $\|E\|$ and $\|Eg\|_\infty$ satisfy, up to constants, the same high probability bounds as obtained through Lemma~\ref{lemma:Wignermatrices} for $\frac1{\sqrt{n}}\left(W - \ddiag(W)\right)$ and so the proof of Theorem~\ref{lem:Z2Synch:exactrecovery} applies.
\end{proof}

% \begin{remark}
%  \TODO{Note that if it is the global solution of the SDP, it needs to be the ground truth~(\cite{?})}
% \end{remark}



\section{Proof of the main results}\label{section:mainresults}

% \TODO{Is there any potential issue in the proofs coming from some non smoothness of the parametrization?} \TODO{uniformize ntoation : cost matrices still different in all sections ...}

This section contains the main technical content of the paper. In particular, it provides guarantees for the Burer--Monteiro approach based solely on deterministic properties of the matrices involved. In this whole section, the cost matrix which appears in~\eqref{rank2} and~\eqref{rank2BM} is denoted by $A$, to mark that the analysis applies both the $\mathbb{Z}_2$-synchronization and to community detection, where the cost matrices are denoted respectively by $Y$ and $A^\sharp$.

%Recall the problem of interest~\eqref{rank2BM} where $Q = [x \ y] \in \reals^{n\times 2}$:
%
%
%
%\begin{equation} 
%\max_{x,y\in\Rn}\quad \langle A, xx^T+yy^T \rangle \quad \text{subject to} \quad x_{i}^2+y_i^2 = 1, \quad i = 1,2,\ldots, n,
%\end{equation}
%and the equivalent re-parametrization by taking $x_i = \cos\theta_i$, $y_i = \sin\theta_i$:
%\begin{equation} \label{reparam}
%\max_{\theta\in\Rn} \quad \cos \theta^T A \cos\theta + \sin\theta^T A \sin\theta, \quad  \theta_i \in \mathbb{R}, \quad i = 1,2,\ldots, n
%\end{equation}
%where we use the notation $\cos \theta = (\cos\theta_1, \ldots, \cos\theta_n)$ and $\sin \theta = (\sin\theta_1, \ldots, \sin\theta_n)$. 



\subsection{Partial Recovery}

%Consider maximizing $ x^T A x + y^T A y$ subject to $x_i^2 + y_i^2 =1$. By parametrizing $x_i = \cos\theta_i$, $y_i = \sin\theta_i$ this is equivalent to maximizing $ cos \theta^T A \cos\theta + \sin\theta^T A sin \theta$ for $ \theta_i \in [0,2\pi]$.

%\textcolor{red}{Note, we need to use one of BM lemma's equating local minima of the factorization relaxation and the rank-constrained problem}

Our first goal is to characterize the local optima of problem~\eqref{rank2}. In particular, we mean to show that they necessarily reveal a lot of information about the ground truth (the planted solution), under some conditions on the noise. To this end, we start by deriving the first- and second-order necessary optimality conditions of~\eqref{rank2BM}.

\begin{lemma} \label{lem:optimality}
	If $Q$ is a local maximum of \eqref{rank2BM} with cost matrix $A$, it satisfies first-order necessary optimality conditions,
	\begin{align}
	\left( \ddiag(AQQ^T) - A \right)Q = 0,
	\label{eq:firstorderconditionmatrix}
	\end{align}
	and second-order necessary optimality conditions,
	\begin{align}
	\ddiag(AQQ^T) - A\circ(QQ^T) & \succeq 0,
	\label{eq:secondorderconditionmatrix}
	\end{align}
	where $\ddiag \colon \Rnn \to \Rnn$ sets all off-diagonal entries to zero. Letting $Q = [x \ y]$, condition~\eqref{eq:firstorderconditionmatrix} also implies
	\begin{align}\label{eq:firstordercondition}
		Ax \circ y = Ay \circ x.
	\end{align}
	(In fact, they are equivalent.)
\end{lemma}
\begin{proof}
	Problem~\eqref{rank2BM} is a smooth optimization problem on a smooth manifold.
	The necessary optimality conditions correspond respectively to requiring the (projected) gradient of the cost function to vanish and the (projected) Hessian of the cost function to have appropriate curvature, see~\citep[eqs.(3.37,5.15)]{AMS08}. A derivation of exactly these conditions is done in full in~\citep{journee2010low,boumal2015staircase}, among others.
	
	We now prove~\eqref{eq:firstorderconditionmatrix} implies~\eqref{eq:firstordercondition}. Note that
	\begin{align*}
		(AQQ^T)_{ii} & = ([Ax \ Ay] Q^T)_{ii} = (Ax)_i x_i + (Ay)_i y_i.
	\end{align*}
	Hence, for all $1 \leq i \leq n$, considering the first and second columns of condition~\eqref{eq:firstorderconditionmatrix} separately, we get
	\begin{align*}
		((Ax)_i x_i + (Ay)_i)x_i & = (Ax)_i, \textrm{ and }\\
		((Ax)_i x_i + (Ay)_i)y_i & = (Ay)_i.
	\end{align*}
	Multiply the first equations by $y_i$ and the second equations by $x_i$, then subtract. It follows for all $i$:
	\begin{align*}
		0 = (Ax)_i y_i - (Ay)_i x_i,
	\end{align*}
	which can be compactly written as $Ax \circ y = Ay \circ x$.
\end{proof}

% Under this parametrization we have
%\[
%\nabla_{\theta} \left(cos \theta^T A \cos\theta + \sin\theta^T A sin \theta \right) = -2(A cos \theta) \circ sin \theta + 2( A \sin\theta) \circ cos \theta
%\]
%and
%\begin{align*}
%\nabla_{\theta\theta} &= 2( A \circ \sin\theta \sin\theta^T) + 2 (A \circ \cos\theta \cos\theta^T) - 2\operatorname{diag}( (A \cos\theta) \circ \cos\theta + (A \sin\theta) \circ \sin\theta)\\
%&= 2A \circ( M(\theta) ) - 2\operatorname{diag}( (A \cos\theta) \circ \cos\theta + (A \sin\theta) \circ \sin\theta)
%\end{align*}
%where $$M(\theta) = \cos(\theta)\cos(\theta)^T + \sin(\theta)\sin(\theta)^T$$ (which is the optimization variable in the rank-2 reduced problem). Thus the conditions for a local max are
%\[
%(A cos \theta) \circ sin \theta = ( A \sin\theta) \circ cos \theta
%\]
%and
%\[
%A \circ M(\theta) - \operatorname{diag}( (A \cos\theta) \circ \cos\theta + (A \sin\theta) \circ \sin\theta) \preceq 0
%\]

\begin{lemma}\label{lemma:Xcorrelates11T}
	Let $A = zz^T + \sigma \Delta$ for some planted signal $z \in \{ \pm 1 \}^n$, where $\Delta=\Delta^T$, $\frac1nSDP(\Delta) \leq \gamma \sqrt{n}$ and $\sigma = c\sqrt{n}$ for some $\gamma, c \geq 0$. Then, for any $Q$ satisfying the second-order condition~\eqref{eq:secondorderconditionmatrix} and corresponding $X = QQ^T$, we have
	%for $M(\theta) = \cos \theta \cos \theta^T + \sin \theta \sin\theta^T$,
	\[
	1 \geq \frac{\langle zz^T, X \rangle}{n^2} \geq \frac{1}{2} - 2\gamma c.
	%\left(\frac{1 - 4\gamma c - (4\gamma c)^2}{1-4\gamma c}\right)^2 - 2\gamma c \geq 0, \quad \textrm{ if } \quad  8\gamma c \leq \sqrt{5}-1.
	\]
	This is true in particular for all local optima, thus showing nontrivial correlation of all of those with the planted solution as soon as $\gamma c < \frac{1}{4}$. Note also that $\frac1nSDP(\Delta) \leq \|\Delta\|$ (see~\eqref{SDP:eq:inequality}).
\end{lemma}
\begin{proof}
	For any two positive semidefinite matrices $X,Y$, it holds that $\langle X, Y \rangle \geq 0$ and that $X \circ Y \succeq 0$ (Schur's product theorem). In particular, taking the inner product of the second-order optimality condition~\eqref{eq:secondorderconditionmatrix} with $X \circ (zz^T) \succeq 0$ on both sides yields the inequality
	\begin{align*}
		\inner{\ddiag(AX)}{X \circ (zz^T)} \geq \inner{A\circ X}{X \circ (zz^T)}.
	\end{align*}
	Since $\diag(X \circ (zz^T)) = \1$, the left hand side evaluates to $\trace(AX)$, so
	\begin{align*}
	\inner{A}{X} \geq \inner{A}{X \circ X \circ (zz^T)}.
	\end{align*}
	Using $A = zz^T + \sigma \Delta$ gives
	\begin{align*}
		\inner{zz^T}{X} + \sigma \inner{\Delta}{X} \geq \inner{(zz^T)\circ(zz^T)}{X \circ X} + \sigma \inner{\Delta}{X \circ X \circ (zz^T)}.
	\end{align*}
	Notice that $$\inner{(zz^T)\circ(zz^T)}{X \circ X} = \inner{\1\1^T}{X\circ X} = \inner{X}{X} = \|X\|_F^2.$$ Furthermore, both $X$ and $X \circ X \circ zz^T$ are feasible for~\eqref{SDP:fullrank}. Thus, by definition~\eqref{def:SDP:eq} and property~\eqref{SDP:eq:inequality},
	\begin{align}
	\langle zz^T, X \rangle &\geq \|X\|_F^2 - 2\sigma\operatorname{SDP}(\Delta) \geq \|X\|_F^2 - 2\gamma c n^2.
	\label{eq:oneonetransposeX}
	\end{align}
%	where we used H\"older's inequality for matrix norms and the fact that the nuclear norms are fixed: $ \| X \|_* = \trace(X) = \|X \circ X\|_* = \trace(X \circ X) = n$. In the latter, we used Schur's product theorem which implies that $X \succeq 0 \implies X\circ X \succeq 0$. \TODO{Should we talk here about the fact that the inner products with $\Delta$ could be bounded with the SDP value? Not sure that adds to the discussion here; omit? or footnote?}
	%
	%Note that since $\text{Tr}(M(\theta)) = n$ and $M(\theta)$ is psd of rank 2, we get $\|M(\theta)\|_F^2 \geq \frac{n^2}{2}$. Thus the 2nd order condition implies
	%\begin{align*}
	%\cos(\theta)^T A \cos(\theta) + \sin(\theta)^T A \sin(\theta) &\geq \frac{n^2}{2} + \sigma \langle \Delta, M(\theta)\circ M(\theta) \rangle \\
	% &\geq \frac{n^2}{2} - \sigma \|\Delta\| \| M(\theta) \circ M(\theta) \|_*\\
	% & = \frac{n^2}{2} - \sigma n \|\Delta\|
	%\end{align*}
	%
	%Using the definition of $A = 11^T + \sigma \Delta$ we have
	%\[
	%\langle 1, \cos(\theta) \rangle^2 + \langle 1, \sin(\theta) \rangle^2 \geq \frac{n^2}{2} - \sigma n \|\Delta\| - (\|\cos(\theta)\|_2^2 + \|\sin(\theta)\|_2^2) \sigma \|\Delta\| = \frac{n^2}{2} - 2\sigma n \|\Delta\|
	%\]
	%
	%Also note that the entries of $M(\theta)\circ M(\theta)$ are non-negative, and in fact the set of matrices of this form are contained in the set of rank-2 psd matrices with non-negative entries. A random \Deltaigner matrix is likely to have small inner product with any matrix in this set (for instance with high probability any rank-1 matrix in such a set cannot be an eigenvector of \Delta). Thus, if we take $\sigma$ small enough, we should get that 
	%\[
	%\cos(\theta)^T A \cos(\theta) + \sin(\theta)^T A \sin(\theta) = \Omega(n^2)
	%\]
%	Using $\|\Delta\| \leq \gamma\sqrt{n}$ and $\sigma = c\sqrt{n}$, we further have:
%	\begin{equation}
%	\langle 11^T, X \rangle \geq \|X\|_F^2 - 2\gamma cn^2.
%	\label{eq:oneonetransposeX}
%	\end{equation} 
	Since $X \succeq 0$, $\rank(X) = 2$ and $\trace(X) = n$, it holds that $ \frac{n^2}{2} \leq \|X\|_F^2 \leq n^2$, so that 
	\begin{align*}
	\langle zz^T, X \rangle \geq n^2(1/2 - 2\gamma c),
	\end{align*}
	concluding the proof.
\end{proof}
We note that, in the above lemma, if $Q$ satisfies second-order conditions only approximately so that $\ddiag(AX) - A\circ C \succeq -\epsilon I_n$ with $X=QQ\transpose$, then the proof still yields $\frac{\inner{zz\transpose}{X}}{n^2} \geq \frac{1}{2} - 2\gamma c - \frac{\epsilon}{n}$.
% The above lemma still works if $X$ has rank more than 2, but you get a $1/p$ instead of $1/2$ I believe.
\begin{lemma}\label{lemma:Qcorrelates1}
	(Continued from Lemma~\ref{lemma:Xcorrelates11T}.) If furthermore $\|\Delta\|\leq \gamma\sqrt{n}$ and $Q$ also satisfies the first-order condition~\eqref{eq:firstorderconditionmatrix}, then
	\[
	1 \geq \frac{\| Q^Tz \|_2}{n} \geq 1 - 8\gamma c,
	\]
	which, in particular, establishes arbitrarily strong correlation between the planted solution $z$ and any local maximum $Q$, as long as $\gamma c$ is sufficiently small.
\end{lemma}
\begin{proof}
	To prove this stronger claim, we need to show that $\|X\|_F$ is arbitrarily close to $n$ for sufficiently small $\gamma c$ (previously, we only had $\|X\|_F \geq n/\sqrt{2}$). Notice that the closer $X$ is to a rank 1 matrix, the closer its Frobenius norm is to $n$, which is in line with our endeavor. In order to get there, we will improve on the previous lemma by incorporating the first-order optimality conditions (which we did not use in the previous lemma).
	
	For ease of notation, we state the proof for $z = \1$. For the general case, apply the change of variable $A \mapsto \diag(z)A\diag(z)$, which does not require changing the assumptions about $\Delta$.
	
	The conditions on $Q$ are invariant under right-orthogonal transformation. That is, we may freely replace $Q$ by $QR$, where $R$ is an arbitrary $2\times 2$ orthogonal matrix. This allows to assume, without loss of generality, that $Q = [x \ y]$ with $\inner{x}{y} = 0$ (simply take the thin SVD of $Q = U\Sigma V^T$ and pick $R = V$). Expand the first-order condition~\eqref{eq:firstordercondition} to get
	\[
	((\1\1^T +\sigma \Delta)x)\circ y = ((\1\1^T +\sigma \Delta)y)\circ x.
	\]
	Thus,
	\[
	\langle \1,x\rangle y - \langle \1, y \rangle x = \sigma (\Delta y)\circ x - \sigma (\Delta x)\circ y.
	\]
	By taking the squared norm of both sides of the previous equation and using $\inner{x}{y}=0$, we have 
	\begin{align*}
	&\langle \1,x\rangle^2  \|y\|_2^2 + \langle \1, y \rangle^2 \| x \|_2^2 \\
	&\leq \sigma^2 \left( \|(\Delta y)\circ x\|_2 + \| (\Delta x)\circ y\|_2 \right)^2 \\
	& \leq \sigma^2 \left( \| \Delta y\|_2 + \|\Delta x\|_2 \right)^2\\
	&\leq \sigma^2 \left( 2\|\Delta\| \sqrt{n} \right)^2 \\
	&\leq c^2 n(2\gamma n)^2.
	\end{align*}
	Now, notice that $\langle \1, x\rangle^2 + \langle \1, y \rangle^2 = \langle \1\1^T, X \rangle$ and the previous lemma imply that either $\langle \1,x\rangle^2$ or $\langle \1,y\rangle^2$ (or both) is larger than or equal to $n^2(1/4 - \gamma c)$. Without loss of generality, assume it is the former:
	\[
	\langle \1,x\rangle^2 \geq n^2(1/4 - \gamma c).
	\]
	%Assuming the right hand side is positive, namely, $4\gamma c < 1$, 
	We may combine the above inequalities to get
	\[
	n^2(1/4 - \gamma c) \|y\|_2^2 \leq \langle 1,x\rangle^2 \|y\|_2^2 \leq (2\gamma c)^2 n^3,
	\]
	or equivalently
	\[
	\|y\|_2^2 \leq 4(2\gamma c)^2 n + 4\gamma c \|y\|_2^2.
	\]
	We aim to show that $\|x\|^2$ is close to $n$. Use $\|x\|_2^2+\|y\|_2^2 = n$:
	\[
	n - \|x\|_2^2 \leq (4\gamma c)^2 n + 4\gamma c n - 4\gamma c \|x\|_2^2,
	\]
	or equivalently
	\[
	(1-4\gamma c) \|x\|_2^2 \geq \left(1 - 4\gamma c - (4\gamma c)^2\right)n.
	\]
	Assuming $1 - 4\gamma c > 0$, we may divide through:
	\[
	\|x\|_2^2 \geq \frac{1 - 4\gamma c - (4\gamma c)^2}{1-4\gamma c}n.
	\]
	As targeted, if $4\gamma c$ is small enough, the right hand side gets arbitrarily close to $n$. In particular, the inequality is only informative if it is nonnegative, which requires $4\gamma c \leq \frac{\sqrt{5}-1}{2}$. By orthogonality of $x$ and $y$, we have $$\|X\|_F^2 = \|x\|_2^4 + \|y\|_2^4 \geq \|x\|_2^4,$$ so that, going back to~\eqref{eq:oneonetransposeX} and under the assumption on $4\gamma c$,
	$$
	\inner{\1\1^T}{X} \geq \|X\|_F^2 - 2\gamma c n^2 \geq \left(\left(\frac{1 - 4\gamma c - (4\gamma c)^2}{1-4\gamma c}\right)^2 - 2\gamma c\right) n^2.
	$$
	In the considered interval, the right hand side is nonnegative if and only if $0 \leq 8\gamma c \leq 1$ (a stronger condition than before). In that interval, we may extract the square root to characterize $Q^T\1$:
	\begin{align*}
	n \geq \|Q^T\1\|_2 = \sqrt{\inner{\1\1^T}{X}} \geq n\sqrt{\left(\frac{1 - 4\gamma c - (4\gamma c)^2}{1-4\gamma c}\right)^2 - 2\gamma c} \geq (1-8\gamma c)n.
	\end{align*}
	The last inequality holds by concavity of the square root term in that interval. We proved the inequality holds for $8\gamma c \leq 1$. It trivially holds otherwise as well. In passing, we note that by restricting $c$ to a smaller interval, the bound can be improved arbitrarily close to $(1-\gamma c)n$ (because less is lost in lowerbounding the concave function by an affine function).
	%
	%
	%
	%Using $\|\sin \theta\|_2^2 \leq n$, we find that
	%\[
	%\|\sin(\theta)\|_2^2 \leq 4(2\gamma c)^2 n + 4\gamma c \|\sin \theta\|_2^2 = 4\gamma c(4\gamma c + 1)n,
	%\]
	%and thus, using $\|\cos\theta\|_2^2+\|\sin\theta\|_2^2 = n$,
	%\[
	%\|\cos(\theta)\|_2^2 \geq \left( 1 - 4\gamma c(4\gamma c + 1) \right)n.
	%\]
	%Using orthogonality of $\cos\theta$ and $\sin\theta$ again,
	%\[
	%\|X\|_F^2 = \|\cos(\theta)\|_2^4 + \|\sin(\theta)\|_2^4 \geq \|\cos(\theta)\|_2^4 \geq  \left(1 - 8\gamma c (4\gamma c + 1)\right) n^2.
	%\]
	%Going back to~\eqref{eq:oneonetransposeX}, we obtain
	%\begin{align*}
	%n^2 \geq \langle 1, \cos(\theta) \rangle^2 + \langle 1, \sin(\theta)\rangle^2 = \langle 11^T, X \rangle & \geq \left(1 - 8\gamma c (4\gamma c + 1) - 2\gamma c \right)n^2 \\
	%& = \left(1 - 10\gamma c - 32(\gamma c)^2 \right)n^2.
	%\end{align*}
	%As targeted, making $c$ sufficiently small now allows to tighten the bounds arbitrarily strongly. It remains to turn these bounds into a statement about $Q$, where $X = QQ^T$. Assuming $\gamma c \leq 1/20$ (conservatively, to make $1 - 10\gamma c - 32(\gamma c)^2 > 0$), we have
	%\begin{align*}
	%	n \geq \|Q^T1\|_2 = \sqrt{\inner{11^T}{X}} \geq \sqrt{1-10\gamma c - 32(\gamma c)^2} n.
	%\end{align*}
	%The square root in the right hand side is a concave function of $\gamma c$. Thus, since we restrict our attention to $\gamma c \leq 1/20$, we may lowerbound it by a linear function:
	%\begin{align}
	%	n \geq \|Q^T1\|_2 \geq (1-12\gamma c)n = \left( 1 - 12\gamma \frac{\sigma}{\sqrt{n}} \right)n,
	%\end{align}
	%which concludes the proof.
	%
	%\TODO{throw away the rest when safe}
	%
	%Let $\alpha = \frac{36c^2}{1/4-3c}$ and $\beta = 6c$. Then we have
	%\[
	%\| 1^T Q \|_2  = \sqrt{ \trace(1^T QQ^T1)} = \sqrt{\trace( 11^T QQ^T) } = \sqrt{\langle 11^T, X \rangle } \geq n - ( 1- \sqrt{(1-\alpha)^2 -\beta})n.
	%\]
	%Using the inequality
	%\[
	%1- \sqrt{(1-\alpha)^2 -\beta} \leq \frac{2\alpha+\beta}{1+ \sqrt{(1-\alpha)^2 -\beta}},
	%\]
	%which holds for $\alpha \geq 0, \beta \geq 0$ sufficiently small, we get
	%\[
	%\| 1^T Y \|_2 \geq n\left(1- \frac{2\alpha+\beta}{1+ \sqrt{(1-\alpha)^2 -\beta}}  \right).
	%\]
	%Taking $c \ll 1$ to be a small enough fixed positive constant we finally get
	%\[
	%\| 1^T Y \|_2 \geq n\left( 1- 4 \frac{\sigma}{\sqrt{n}} \right),
	%\]
	%which concludes the proof.
\end{proof}
Thus, by taking $\sigma = c\sqrt{n}$ for constant $c > 0$ small enough, we have that any local maximizer $Q$ (in fact, any point satisfying both first- and second-order necessary optimality conditions) correlates very strongly with $z$.



\subsection{Exact recovery}

To establish exact recovery, we only need to show that all second-order critical points of~\eqref{rank2BM} have rank 1. Indeed, it is well-known that rank deficient second-order critical points $Q$ are global optima.
\begin{lemma}\label{lem:rankdeficientglobalopt}
	If $Q$ satisfies the second-order necessary optimality condition for~\eqref{rank2BM} and it is rank deficient (i.e., $\rank(Q) = 1$), then $Q$ is a global optimum.
\end{lemma}
\begin{proof}
	Given $Q$ feasible for~\eqref{rank2BM} with cost matrix $A$, observe that if
	\begin{align}
		S & = S(Q) = \ddiag(AQQ^T) - A
		\label{eq:S}
	\end{align}
	is positive semidefinite, then $Q$ is a global optimum. Indeed, for any feasible contender $\tilde Q$, we have (using $\diag(\tilde Q \tilde Q^T) = \1$)
	\begin{align*}
		0 \leq \innersmall{S}{\tilde Q \tilde Q^T} = \trace(AQQ^T) - \trace(A\tilde Q \tilde Q^T),
	\end{align*}
	thus showing that $\trace(Q^TAQ)$ is maximal over the search space. The matrix $S$ features prominently in~\citep{burer2002rank,burer2005local,journee2010low,wen2013orthogonality,bandeira2014tightness,boumal2015staircase,boumal2016nonconvexphase}, understandably given its strong properties as an optimality certificate. In this paper, we rely less on it in favor of different arguments.
	
	If $\rank(Q) = 1$, then $QQ^T = qq^T$ for some $q \in \{ \pm 1 \}^n$. We show that if $Q$ further satisfies the second-order condition~\eqref{eq:secondorderconditionmatrix}, then $S(Q)$ is positive semidefinite, implying optimality of $Q$. For all $u\in\Rn$, we have (using $(QQ^T)_{ij}^2 = 1$ for all $i,j$)
	\begin{align*}
		u^TSu & = \inner{\ddiag(AQQ^T) - A}{uu^T} \\
			  & = \inner{\ddiag(AQQ^T) \circ (QQ^T) - A \circ (QQ^T)}{(uu^T) \circ (QQ^T)} \\
			  & = \inner{\ddiag(AQQ^T) - A \circ (QQ^T)}{(u\circ q)(u\circ q)^T} \geq 0,
	\end{align*}
	where the inequality follows from~\eqref{eq:secondorderconditionmatrix}. This holds for all $u$, thus $S\succeq 0$.
\end{proof} 
%\TODO{Afonso, comment out this remark after you read it: the proof above is a bit different from the traditional ones; I think a messy version of this could show that a point where the Hessian is approximately psd and the matrix is approximately rank deficient (so, all up to numerical tolerances) has to be approximately optimal, and we would get an actual bound from that. It wouldn't look nice, but it's there, and I don't think it was before.}
We further show in the present context another well-known result, namely that if the perturbation $\sigma \Delta$ has good properties, then the SDP has a unique solution corresponding to the ground truth. In the Wigner setting, this lemma may be improved somewhat, see~\citep[\S2.1]{bandeira2014tightness}.
\begin{lemma}\label{lem:uniqueglobaloptexact}
	Let $A = zz^T + \sigma \Delta$ for some planted solution $z \in \{\pm 1\}^n$, with $\sigma = c\sqrt{n}$, $\Delta = \Delta^T$, $\diag(\Delta) = 0$, $\|\Delta\| \leq \gamma \sqrt{n}$ and $\|\Delta z\|_\infty \leq \gamma \sqrt{n \log n}$, for some $\gamma, c \geq 0$. If
	\begin{align*}
		\gamma c < \frac{1}{1+\sqrt{\log n}},
	\end{align*}
	then the unique global optimum of~\eqref{rank2} with cost matrix $A$ is $X = zz^T$, so that all global optima $Q$ of~\eqref{rank2BM} are of the form $Q = [z \ 0]R$, where $R$ is a $2\times 2$ orthogonal matrix.
\end{lemma}
\begin{proof}
	We show that $S = S(z) = \ddiag(Azz^T) - A$~\eqref{eq:S} is positive semidefinite with rank $n-1$, which by the argument in the proof of Lemma~\ref{lem:rankdeficientglobalopt} implies optimality of $zz^T$. Uniqueness follows from strict complementarity ($\rank(zz^T) + \rank(S(z)) = n$).
	Indeed: let $X$ be any optimum of~\eqref{rank2}; then, $\inner{S}{X} = \inner{A}{zz^T} - \inner{A}{X} = 0$. Since $S \succeq 0$ and $X \succeq 0$, $\inner{S}{X} = 0$ implies $S X = 0$. Thus, $\spann(X) \subset \ker S$. But $\ker S = \spann(z)$ since the kernel has dimension~1 and $Sz = 0$. Adding that $X \succeq 0$ and $\diag(X)=\1$, it comes that $X = zz\transpose$.
	%; see a proof in~\citep[Thm.~3.4]{boumal2015staircase}.
	
	Using $Sz = 0$, it remains to show that for all $u \neq 0$ such that $z^Tu = 0$, $u^TSu > 0$. We have:
	\begin{align*}
		u^TSu & = u^T\left( nI_n - zz^T + \sigma \left[ \ddiag(\Delta zz^T - \Delta) \right] \right)u \\
			& = n\|u\|_2^2 + \sigma \left[ u^T\ddiag(\Delta zz^T)u - u^T\Delta u \right] \\
			& \geq \|u\|_2^2 \left( n - \sigma \|\Delta z\|_\infty - \|\Delta \| \right) \\
			& \geq n \|u\|_2^2 \left( 1 - \gamma c \sqrt{\log n} - \gamma c  \right).
	\end{align*}
	This is indeed positive under the prescribed condition.
\end{proof}
We go through a few technical lemmas to establish the rank-one property of second-order critical points, under some conditions on the perturbation.
\begin{lemma}\label{lem: MaxCutOptimality}
	Let $QQ^T$ be any feasible point of \eqref{rank2} satisfying first- and second-order necessary optimality conditions, with $Q = [x\ y] \in \mathbb{R}^{n\times 2}$. Let $A_i$ be the $i$th row of $A$. If $\diag(A) \geq 0$, then
	\[
	\forall i, \quad \diag(AQQ^T)_i = \langle A_i, x \rangle x_i + \langle A_i, y \rangle y_i = \sqrt{ \langle A_i, x\rangle ^2 + \langle A_i, y \rangle^2 } = \|e_i^TAQ\|_2.
	\]
\end{lemma}
\begin{proof}
	Row-wise, the first-order condition \eqref{eq:firstorderconditionmatrix} reads
	\[
	\left [ \langle A_i, x \rangle, \langle A_i, y \rangle \right] = \left( \langle A_i, x\rangle x_i + \langle A_i, y \rangle y_i \right) \left[x_i, y_i \right].
	\]
	Taking the $L_2$ norm of both sides and using $x_i^2 + y_i^2 = 1$, we get
	\[
	\sqrt{ \langle A_i, x\rangle ^2 + \langle A_i, y \rangle^2 } = \left| \langle A_i, x \rangle x_i + \langle A_i, y \rangle y_i \right|.
	\]
	The second-order optimality condition~\eqref{eq:secondorderconditionmatrix} implies that (simply considering the diagonal of the positive semidefinite matrix, which must be nonnegative)
	\[
	\diag(AQQ^T) \geq \diag(A).
	\]
	That is,
	\[
	\langle A_i, x \rangle x_i + \langle A_i, y \rangle y_i  \geq A_{ii}.
	\]
	Since we assume $\diag(A) \geq 0$, the two results combine into:
	\begin{align}
	\langle A_i, x \rangle x_i + \langle A_i, y \rangle y_i = \left| \langle A_i, x \rangle x_i + \langle A_i, y \rangle y_i \right| = \sqrt{ \langle A_i, x\rangle ^2 + \langle A_i, y \rangle^2 },
	\end{align}
	concluding the proof.
\end{proof}
%
%\textcolor{red}{
%\begin{lemma}\label{lem: maxrowWY}
% For $Y$ satisfying second-order optimality and $\sigma = o(\sqrt{n})$ conditions, we have $\max_i \|e_i^TWY\|_2 \lesssim \sqrt{n} + \left(\frac{\sigma}{\sqrt{n}}\right)^2 n$.
%\end{lemma}
%\begin{proof}
%	See handwritten notes on page 6 and 8 of ``ForNicolasAndVlad.pdf'' and use Lemma~\ref{lemma:correlationof1withY_0}.
%\end{proof}
%}
%
The following lemma bounds the largest row in $\Delta Q$. For a Wigner setting, this result is likely suboptimal, and is the bottleneck in our analysis. This is the same bottleneck that arose in both~\citep{bandeira2014tightness} and~\citep{boumal2016nonconvexphase} for the phase synchronization problem.
\begin{lemma}\label{lem:maxrowWY2}
	Let $QQ^T$ be any feasible point of \eqref{rank2} satisfying first- and second-order necessary optimality conditions with cost matrix $A = zz^T + \sigma \Delta $, $\sigma = c\sqrt{n}$, where $\Delta  = \Delta^T$ satisfies $\diag(\Delta ) = 0$, $\|\Delta \| \leq \gamma \sqrt{n}$ and $\|\Delta z\|_\infty \leq \gamma \sqrt{n \log n}$, for some $\gamma, c \geq 0$. Then, we have
	$$
	\max_i \|e_i^T\Delta Q\|_2 \leq \gamma \sqrt{n} \left( \sqrt{\log n}  + 4\sqrt{\gamma cn}\right).
	$$
\end{lemma}
\begin{proof}
	Once more, we write the proof for $z=\1$, without loss of generality.
	Let $P_\1 = \frac{1}{n} \1\1\transpose$ be the orthogonal projector to $\text{span}(\1)$ where $\1 \in\Rn$ and, analogously, let $P_{\1^\perp} = I-P_\1$ be the projector to $\text{span}(\1)^\perp$. Writing $w_i$ for the $i$th column of $\Delta$ and letting $Q = [x\ y]$, we have for all $i$:
	\begin{align*}
	\left\|e_i^T \Delta  Q \right\|_2 &= \left\| w_i^T (P_{\1} + P_{\1^\perp})Q\right\|_2\\
	&\leq \left\|w_i^T P_\1 (Q) \|_2 + \|w_i^T P_{\1^\perp} (Q) \right\|_2\\
	& \leq \frac{1}{n}\left\| w_i^T \1\1^TQ \right\|_2 + \|w_i\|_2 \left\| Q - \frac{1}{n} \1\1^TQ \right\|_F \\
	& \leq \frac{1}{n} \|\Delta \1\|_\infty \|Q^T \1\|_2 + \|w_i\|_2 \left\| Q - \frac{1}{n} \1\1^TQ \right\|_F.
	\end{align*}
	Since,
	%Note that the assumptions of lemma \ref{lemma:Vlad:11TcorrelateswithY} are satisfied with high probability by standard concentration inequalities for Wigner matrices. Thus
	by Lemma~\ref{lemma:Qcorrelates1},
	$\|Q^T \1\|_2 \geq n\left(1 - 8\gamma c\right)$, we may further bound the last term using
	\begin{align*}
	\left\| Q - \frac{1}{n} \1\1^TQ \right\|_F^2 & = \|Q\|_F^2 + \frac{1}{n^2} \|\1 \1^TQ\|_F^2 - \frac{2}{n} \|Q^T \1\|_2^2 \\
	& = n - \frac{1}{n} \|Q^T \1\|_2^2 \\
	& \leq n ( 1 - (1-8\gamma c)^2) \leq 16\gamma c n,
	\end{align*}
	where we used $\|Q\|_F^2 = n$.
	%where we used that $\|e_i^T Y\|_2^2 = 1$ and lemma \ref{lemma:Vlad:11TcorrelateswithY}.
	Combining and using $\|Q^T\1\|_2 \leq n$, we get for all $i$ that
	\begin{align*}
	\left\|e_i^T \Delta  Q \right\|_2 & \leq \gamma \sqrt{n \log n} + \|\Delta \| \sqrt{16\gamma cn} \\
	& \leq \gamma \sqrt{n} \left( \sqrt{\log n}  + 4\sqrt{\gamma cn}\right).
	\end{align*}
	This concludes the proof.
\end{proof}
The next lemma is central to our endeavor: it identifies a regime in which all second-order critical points have rank 1. The first inequality used in this proof is an important step inspired by the proof of~\citep[Thm.~3]{wen2013orthogonality}. Crucially, it is because we inject both first- and second-order optimality conditions (as opposed to only first-order conditions in~\citep{wen2013orthogonality}) that we are able to make a substantial statement via this rank-control argument. Indeed, even for the noiseless case, Theorem 3 in \citep{wen2013orthogonality} does not lead to sufficiency of $p=2$ for SDPLR.
\begin{lemma} \label{lem:master}
	Under the assumptions of Lemma~\ref{lem:maxrowWY2} and using the same notation, if
	\begin{align*}
		\gamma c < \frac{1}{9 + \sqrt{\log n} + 4\sqrt{\gamma c n}},
	\end{align*}
	then $\rank(Q) = 1$.
\end{lemma}
\begin{proof}
	From the first-order condition~\eqref{eq:firstorderconditionmatrix}, we may control the rank of $Q$ as
	\begin{align*}
	\rank(Q) & \leq \nulll\left(\ddiag(AQQ^T)-A\right) \\
	& = n - \rank\left(\ddiag(AQQ^T)-A\right) \\
	& = n - \rank\left(\ddiag(AQQ^T)-\sigma \Delta  - zz^T\right) \\
	& \leq n+1 - \rank\left(\ddiag(AQQ^T)-\sigma \Delta \right).
	\end{align*}
	To ensure $\rank(Q) = 1$, it remains to force $\rank(\ddiag(AQQ^T)-\sigma \Delta ) = n$. Since the first matrix is diagonal, this is the case in particular if
	\begin{align*}
	\min_i \diag(AQQ^T)_i > \sigma \|\Delta \|.
	\end{align*}
	This can be controlled by Lemmas~\ref{lem: MaxCutOptimality} and~\ref{lem:maxrowWY2}:
	\begin{align*}
	\min_i \|e_i^T AQ\|_2 - \sigma \|\Delta \| & = \min_i \|e_i^T zz^T Q + \sigma \cdot e_i^T \Delta  Q\|_2 - \sigma \|\Delta \| \\
	& \geq \|Q^Tz\|_2 - \sigma \cdot \max_i \|e_i^T \Delta Q\|_2 - \gamma c n \\
	& \geq n - 9\gamma c n - \gamma c n \left( \sqrt{\log n}  + 4\sqrt{\gamma cn}\right).
	\end{align*}
	Forcing the latter to be positive (as with the condition in this lemma's statement) is sufficient to imply $\rank(Q) = 1$.
	%
	%\TODO{STOP}
	%	
	%By lemma \ref{lem: MaxCutOptimality}, we have
	%\begin{align}
	%	 \langle A_i, x \rangle x_i + \langle A_i, y \rangle y_i  \leq -1
	%\end{align}
	%so that $| \langle A_i, x \rangle x_i + \langle A_i, y \rangle y_i| = - \left(\langle A_i, x \rangle x_i + \langle A_i, y \rangle y_i\right)$. Thus, again using lemma \ref{lem: MaxCutOptimality} we get
	%\begin{align}
	%	- \left(\langle A_i, x \rangle x_i + \langle A_i, y \rangle y_i\right) = \| A_i^T Y \|_2 = \| 1\transpose Y - \sigma e_i^T (WY) \|_2 \geq  \|1\transpose Y\|_2 - \sigma \max_i \|e_i^TWY)\|_2.
	%	\label{eq:CYYk}
	%\end{align}
	%Therefore,
	%\[
	%-\operatorname{diag}(AYY^T) \succeq \gamma n I_{n \times n}
	%\]
	%
	%By lemma \ref{lem: MaxCutOptimality} again, we have $(A - \operatorname{diag}(AYY^T))Y = 0$, so that
	%\begin{align*}
	%	\rank Y \leq \nulll(A - \operatorname{diag}(AYY^T)) & = n - \rank(A - \operatorname{diag}(AYY^T)) \\
	%	 & = n - \rank\left( \sigma W + 11\transpose - \ddiag(AYY\transpose) \right) \\
	%	 & \leq n+1 - \rank(\sigma W - \ddiag(AYY\transpose)).
	%\end{align*}
	%Since $\sigma W \prec \gamma n I_{n \times n}$ and $-\ddiag(AYY\transpose) \succeq \gamma n I_n$, we have
	%\begin{align}
	%	\rank(\sigma W - \ddiag(AYY\transpose)) = n.
	%\end{align}
	%and hence
	%\[
	%\rank(Y) = 1
	%\]
	%
\end{proof}


\begin{theorem} \label{thm:maxcutrktwo}
	Let $A = zz^T + \sigma \Delta $ for some planted solution $z \in \{\pm 1\}^n$, with $\sigma = c\sqrt{n}$, $\Delta  = \Delta^T$, $\diag(\Delta ) = 0$, $\|\Delta \| \leq \gamma \sqrt{n}$ and $\|\Delta z\|_\infty \leq \gamma \sqrt{n \log n}$, for some $\gamma, c \geq 0$. If
	\begin{align*}
		\gamma c < \frac{1}{9 + \sqrt{\log n} + 4\sqrt{\gamma c n}},
	\end{align*}
	then all second-order critical points $Q$ of~\eqref{rank2BM} with cost matrix $A$ are global optima of rank 1 such that $QQ^T = zz^T$.
	%If the planted solution is $x \in \{\pm 1\}^n$ instead of 1, the theorem holds with $QQ^T = xx^T$ provided the conditions on $W$ and $W1$ hold for $\diag(x)W\diag(x)$ and $\diag(x)W\diag(x)$.
%	
%	\TODO{careful with the change of variable $1 \to x$; should be good though.}
%	
%	\TODO{Former statement: Consider the maximization problem \eqref{rank2BM} and let $A = xx^T + \sigma W$ where $W$ is a standard Wigner random matrix with $\operatorname{diag}(W) = 0$, where $x\in\mathbb{R}^n$ with $x_i^2 =1$. Then there exists an absolute constant $c>0$ such that if $\sigma \leq cn^{1/6}$, then with very high probability every local maximum $Y$ of \eqref{rank2BM} is rank 1, with column span of $Y$ equal to $\text{span}(x)$. %Thus, with high probability maximizing the rank-2 relaxation \ref{rank2} exactly recovers the matrix $xx^T$. 
%	}
	There exists a constant $k$ such that, if $\gamma c \leq k n^{-1/3}$, then the theorem applies.
\end{theorem}
\begin{proof}
	By Lemma~\ref{lem:master}, all second-order critical points $Q$ have rank 1. By Lemma~\ref{lem:rankdeficientglobalopt}, such $Q$'s are thus globally optimal. By Lemma~\ref{lem:uniqueglobaloptexact} (whose conditions are satisfied a fortiori), they all satisfy $QQ^T = zz^T$.
%	By theorem \textcolor{red}{cite BM} we can wlog instead consider any local maximum of problem \eqref{rank2BM}. Since the trace functional and the distribution of a Wigner matrix are invariant to conjugation by rank 1 sign matrices, we can wlog take $x = 1 \in \mathbb{R}^n$. Applying lemmas \ref{lemma:Vlad:11TcorrelateswithY} and \ref{lem: maxrowWY2}, and using the assumption on $\sigma$ we see that the assumptions of lemma \ref{lem: master} are satisfied. That is, since whp $\|W\| \leq 3 \sqrt{n}$ we have by lemma \ref{lemma:Vlad:11TcorrelateswithY}
%	\[
%	\|1^T Y\|_2 \geq \left(1- 4 \frac{\sigma}{\sqrt{n}}\right)n
%	\]
%	and by lemma~\ref{lem:maxrowWY2} that whp
%	\[
%	\max_i \| e_i^T WY \|_2 \leq \frac{\text{polylog}(n)}{\sqrt{n}} \|1^T Y\|_2 + 3 \sqrt{8} n^{3/4} \sigma^{1/2}
%	\]
%	Thus, whp
%	\[
%	\|1^T Y\|_2 - \sigma\max_i \| e_i^T WY \|_2 \geq n\left(1 - 4\frac{\sigma}{\sqrt{n}}\right)\left(1 - \sigma \frac{\text{polylog}(n)}{\sqrt{n}}\right) -  3 \sqrt{8} n^{3/4} \sigma^{3/2}
%	\]
%	By taking $c$ to be small enough in the definition of $\sigma$, and combining with $\|\sigma W\| \ll n$ whp, we thus have that the assumptions of lemma \ref{lem: master} are satisfied. Thus, with high probability every local maximizer is rank 1. By proposition \ref{prop:YrankdeficientKKT2}, $YY^T$ is then optimal for the SDP and by \textcolor{red}{citation} $YY^T = 11^T$. 
\end{proof}


%any local maximum must be highly correlated with the rank 1 matrix $11^T$. Note that if 
%\[
%\langle 11^T, M(\theta) \rangle =  \|M(\theta)\|_F^2
%\]
%then $M(\theta) = 11^T$. Note that we haven't yet used the first order optimality condition. Moreover we have not used the full strength of the 2nd order optimality condition. The rest of the proof could follow by performing a local analysis of $11^T$ to show that there is a radius of convergence or by using the 1st order optimality condition.

%The first order condition encourages $\cos\theta = \sin\theta$, which is what we want to show, and the 2nd order condition encourages $\cos\theta$ and $\sin\theta$ to be a large eigenvector of $A$. Recall that in our case $A = 11^T + \sigma W$ which has a large rank 1 component in the direction $1$. $A = 11^T + \sigma W$ of course had plenty of positive eigenvectors, which is what we want to exploit to break the 2nd order condition. This would be difficult to exploit if $M$ somehow zeroed out those positive eigenvectors, but the structure of $M$ seems to prevent this, as there have to be enough large positive entries of $M$ (as long as we can  take $cos \theta \geq 0$ and $\sin\theta \geq 0$, which may follow from direct manipulation of the objective function). Thus, there should be large sub-blocks of $M$ that are $>= c11^T$ (for vector 1 of appropriate dimension), and since eigenvectors of a Wigner are random, they are bound to be incoherent (non-sparse), which means they would have positive correlation with such sub-block.   

%\TODO{Add a remark explaining that the suboptimality is the same issue as in \citep{bandeira2014tightness,boumal2016nonconvexphase}, formulate conjecture.}

% Acknowledgments---Will not appear in anonymized version
\acks{ASB was supported by NSF grant DMS-1317308. ASB acknowledges Wotao Yin for pointing the author to the relevant reference~\citep{wen2013orthogonality} which helped motivate the start of the investigation in this paper. NB, formerly hosted by the SIERRA group at Inria and ENS, was supported by the ``Fonds Sp\'eciaux de Recherche'' (FSR) at UCLouvain and by the Chaire Havas ``Chaire Eco\-no\-mie et gestion des nouvelles don\-n\'ees'', the ERC Starting Grant SIPA and a Research in Paris grant in Paris. VV acknowledges support from the Office of Naval Research.}



\bibliography{boumal_BM}














 \appendix
% 
% 
% 
% 
% 
 \section{Some technical steps}

 
\subsection{Other needed Lemmas}

%\TODO{Change Theorem to Community Detection, it will be a constant away. Explain the changes}

%\TODO{It is easy to see that for the very first Lemma, where all we need is to get positive correlation, the argument doesn't really depend on $\|W\|$ but on the maximum of $\frac1n \max_{X\succeq 0, X_{ii}=1} \trace(WX)$ and this is known to behave well for community detection, see for example:~\cite{Montanari_SDPdetectionSBM} (Lemma H.2.)}

%\TODO{We also need upper bounds on spectral norm of the ``noise matrix'', these can be obtained from~\cite{Bandeira_NARandomMatrixBound}}

\begin{lemma}\label{lemma:Wignermatrices}
 Let $W$ be a symmetric Wigner matrix whose entries are independent standard gaussian random variables and $z\in\{\pm1\}$ a fixed vector. $W-\ddiag(W)$ is the same matrix with the diagonal elements replaced by zeros. Then, the following holds for any $t\geq 0$:
 \begin{equation}\label{eq:lemma:Wspectralnorm}
  \mathrm{Prob}\left( \|W-\ddiag(W)\| \geq 2\sqrt{n} + t \right) \leq \exp\left( - \frac{t^2}4 \right),
 \end{equation}
 and
 \begin{equation}\label{eq:lemma:Winftynorm}
  \mathrm{Prob}\left( \|(W-\ddiag(W))z\|\infty \geq  \sqrt{2\log {n}+t} \right) \leq \exp\left( - \frac{t}2 \right).
 \end{equation}
\end{lemma}
\begin{proof}
\eqref{eq:lemma:Wspectralnorm} follows from combining $\EE\|W-\ddiag(W)\| \leq \EE\|W\|$ (which follows from Jensen's inequality), the well known fact that $\EE\|W\|\leq 2\sqrt{n}$, and Gaussian Concentration. %, see for example Section 4 in~\cite{bandeira2015lecturenotes}.
\eqref{eq:lemma:Winftynorm} follows from standard tail bounds on gaussian random variables together with a simple union bound.
\end{proof}

Recall the definition of $E$ in~\eqref{eq:def:E}, $E$ is symmetric, its entries are independently distributed, has diagonal zero and the distribution of the entries is: for $i\neq j$ and $g_i=g_j$:
\begin{equation}
\sqrt{\frac{(p+q)n}{2}}E_{ij} = \left\{ \begin{array}{ll}
1 - p  & \text{ with probability }p \\
-p &  \text{ with probability }1-p,
\end{array}\right.
\end{equation}
and, if $g_i\neq g_j$:
\begin{equation}
\sqrt{\frac{(p+q)n}{2}}E_{ij} = \left\{ \begin{array}{ll}
1 - q  & \text{ with probability }q \\
-q &  \text{ with probability }1-q.
\end{array}\right.
\end{equation}


\begin{lemma}\label{lemma:sbm:SDPE}
There exists a constant $C$ such that, with high probability
\[
\mathrm{SDP}(E) \leq 2 + C\frac{\log d}{d^{1/10}},
\]
where $d = \frac{(p+q)n}{2} = \frac{a+b}2$ is the average degree parameter.
\end{lemma}

\begin{proof}
See Lemma~H.2. and equation (259) in~\citep{Montanari_SDPdetectionSBM}. 
\end{proof}

\begin{lemma}\label{lemma:sbm:spectralnormE}
Consider $E$ as defined above. With high probability, there exists a constant $C$ such that
\[
 \|E\| \leq 3 + C\sqrt{\frac{\log n}{\frac{n}2(p+q)}}.
\]
\end{lemma}
\begin{proof}
Set $\tilde{E} = \sqrt{\frac{(p+q)n}{2}}E$. Since the entries of $\tilde{E}$ are independent, we can use Corollary 3.12 in \citep{Bandeira_NARandomMatrixBound} (with, say, $\eps=3$) to obtain
\[
\mathrm{Prob}\left( \|\tilde{E}\| \geq 3\tilde{\sigma} + t \geq  \right) \leq n\exp\left(-\frac{t^2}{c\tilde{\sigma_\ast}^2}\right),
\]
for any $t>0$ and a universal constant $c$. Here
\[
\tilde{\sigma}  = \max_{i}\sqrt{\sum_j\EE E_{ij}^2} \leq \sqrt{\frac{n}2p(1-p) + \frac{n}2p(1-p)}, 
\]
and
\[
\tilde{\sigma_\ast} = \max_{ij}\|E_{ij}\| \leq 1.
\]
This means that, with high probability, there exists a constant $C$ such that
\[
 \|\tilde{E}\| \leq 3\sqrt{\frac{n}2p(1-p) + \frac{n}2q(1-q)} + C\sqrt{\log n} \leq 3\sqrt{\frac{(p+q)n}2} + C\sqrt{\log n},
\]
concluding the proof.
\end{proof}


\begin{lemma}\label{lemma:sbm:infnormE}
 Consider $E$ as defined above. With high probability, there exists a constant $C$ such that
 \[
  \|Eg\|_\infty \leq C\sqrt{\log n} + C\frac{\log n}{\sqrt{\frac{n}2(p+q)}}
 \]
\end{lemma}
\begin{proof}
Set $\tilde{E} = \sqrt{\frac{(p+q)n}{2}}E$. For each $i\in [n]$, we have
\[
 (\tilde{E}g)_i = \sum_{j=1}^{\frac{n}2-1}\xi_j - \sum_{j=1}^{\frac{n}2}\xi'_j,
\]
where $\xi_j$ are independent random variables taking the value $1-p$ with probability $p$ and the value $-p$ with probability $1-p$; $\xi'_j$ are independent random variables (and independent to the random variables $\xi'_{j'}$) taking the value $1-q$ with probability $q$ and the value $-q$ with probability $1-q$. It is easy to see that
\[
 \sum_{j=1}^{\frac{n}2-1}\EE\xi_j^2 - \sum_{j=1}^{\frac{n}2}\EE\left(\xi'_j\right)^2 = \frac12(\frac{n}2-1)p(1-p) + \frac{n}2q(1-q) \leq \frac{n}2(p(1-p)+q(1-q)) ,
\]
and that the summands are almost surely bounded by $1$. This means that we can use Bernstein's inequality to get
\[
 \mathrm{Prob}\left( (\tilde{E}g)_i > t \right) \leq \exp\left( - \frac{\frac12t^2}{\frac{n}2(p(1-p)+q(1-q)) + \frac13t} \right).
\]
A union bound gives
\[
 \mathrm{Prob}\left( \|\tilde{E}g\|_\infty > t \right) \leq 2n\exp\left( - \frac{\frac12t^2}{\frac{n}2(p(1-p)+q(1-q)) + \frac13t} \right),
\]
which means that, with high probability,
\[
 \|\tilde{E}g\|_\infty \lesssim \log n + \sqrt{\frac{n}2(p(1-p)+q(1-q))}\sqrt{\log n} \leq \log n + \sqrt{\frac{n}2(p+q)}\sqrt{\log n},
\]
concluding the proof.
\end{proof}

\begin{remark}
It is worth noting that the quantities $\|E\|$ and $\|Eg\|_\infty$ are tightly connected to the control of the spectrum of $\Gamma_{SBM}$ in~\citep[Definition 4.8]{bandeira2015laplacian}. 
\end{remark}


 
 
% \subsection{Proof of Lemma~\eqref{?}}
% 


%\section{Other SDPs \TODO{remove entire section?}}
%
%\TODO{Write this Section} \TODO{NB: didn't we say we would drop this? Offset it to a journal paper? otherwise, there'll be nothing more to say.}
%
%a) This is a very general method
%
%In general 
%
%
%Semidefinite programs (SDPs) are optimization problems of the form
%\begin{align}
%	f^* = \min_{X\in\Snn} \inner{C}{X} \textrm{ subject to } \calA(X) = b, X \succeq 0,
%	\tag{SDP}
%	\label{eq:SDP}
%\end{align}
%where $\Snn$ is the set of real symmetric $n \times n$ matrices, $\calA \colon \Rnn \to \Rm$ is a linear operator describing $m$ equality constraints, $\inner{A}{B} \triangleq \trace(A\transpose B)$ is a standard inner product over matrices and $X \succeq 0$ constrains $X$ to be positive semidefinite.
%
%
%
%The manifold assumption
%
%A general Theorem (in the companion paper), was somewhat of a folk result
%
%
%b) What about over the complex case? Talk about tightness in our paper and Nicolas paper
%c) Somehow the bottleneck seems to be the same as in the other paper, and it seems to relate to the Open problem in my notes.~\cite{bandeira2015lecturenotes}
%
%
%Sum of Squares stuff
%
%
%\TODO{
%It would be very interesting to say something about Orthogonal Cut~\cite{bandeira2013approximating,boumal2015staircase} or even more general~\cite{bandeira2015nonuniquegames}, mention this in future work.
%}


\end{document}
